\NormalFont\ShortTitle{Ioan}

{\MTB Vestea cea bună conform

\par }{\MT Ioan


\par }\ChapOne{1}{\PP \VerseOne{1}La început era Cuvântul, și Cuvântul era cu Dumnezeu, și Cuvântul era Dumnezeu.

\VS{2}El era la început cu Dumnezeu.

\VS{3}Toate lucrurile au fost făcute prin el. Fără el, nimic din ceea ce a fost făcut nu s-a făcut.

\VS{4}În El era viața, iar viața era lumina oamenilor.

\VS{5}Lumina strălucește în întuneric, iar întunericul nu a biruit-o.


\par }{\PP \VS{6}Și a venit un om trimis de Dumnezeu, care se numea Ioan.

\VS{7}Acesta a venit ca martor, ca să mărturisească despre lumină, pentru ca toți să creadă prin el.

\VS{8}El nu era lumina, ci a fost trimis ca să mărturisească despre lumină.

\VS{9}Lumina adevărată, care luminează pe toată lumea, venea în lume.


\par }{\PP \VS{10}El era în lume, și lumea a fost făcută prin El, și lumea nu L-a recunoscut.

\VS{11}El a venit la ai Săi, și cei ce erau ai Săi nu L-au primit.

\VS{12}Dar tuturor celor care l-au primit, le-a dat dreptul de a deveni copii ai lui Dumnezeu, celor ce cred în numele lui;

\VS{13}care nu s-au născut din sânge, nici din voința cărnii, nici din voința omului, ci din Dumnezeu.


\par }{\PP \VS{14}Cuvântul S-a făcut trup și a locuit printre noi. Am văzut slava Lui, o slavă ca a Fiului unic născut din Tatăl, plin de har și de adevăr.

\VS{15}Ioan a mărturisit despre el. El a strigat, spunând: “Acesta a fost cel despre care am spus: “Cel care vine după mine m-a întrecut pe mine, pentru că a fost înaintea mea”.”

\VS{16}Din plinătatea lui, noi toți am primit har peste har.

\VS{17}Căci Legea a fost dată prin Moise. Harul și adevărul s-au realizat prin Isus Cristos.

\VS{18}Nimeni nu L-a văzut pe Dumnezeu în niciun moment. Fiul unic născut, care este în sânul Tatălui, l-a declarat.


\par }{\PP \VS{19}Iată mărturia lui Ioan, când iudeii au trimis preoți și leviți de la Ierusalim ca să-l întrebe: “Cine ești Tu?”


\par }{\PP \VS{20}El a declarat, și nu a negat, ci a declarat: “Eu nu sunt Hristosul.”


\par }{\PP \VS{21}Și L-au întrebat: “Și atunci, ce? Ești tu Ilie?”

\par }{\PP El a spus: “Nu sunt.”

\par }{\PP “Tu ești profetul?”

\par }{\PP El a răspuns: “Nu.”


\par }{\PP \VS{22}Și I-au zis: “Cine ești Tu? Dă-ne un răspuns pe care să-l ducem înapoi la cei care ne-au trimis. Ce spui despre tine însuți?”


\par }{\PP \VS{23}El a zis: “Eu sunt glasul celui care strigă în pustiu: “Îndreptați calea Domnului”, cum a zis profetul Isaia”.


\par }{\PP \VS{24}Cei care fuseseră trimiși erau din partea fariseilor.

\VS{25}Ei L-au întrebat: “De ce deci botezi, dacă nu ești nici Hristosul, nici Ilie, nici proorocul?”


\par }{\PP \VS{26}Ioan le-a răspuns: “Eu botez în apă, dar în mijlocul vostru este Unul pe care nu-L cunoașteți.

\VS{27}El este cel care vine după mine, care este preferat înaintea mea, căruia nu sunt vrednic să-i dezleg cureaua de la sandale.”

\VS{28}Aceste lucruri se întâmplau în Betania, dincolo de Iordan, unde boteza Ioan.


\par }{\PP \VS{29}A doua zi, a văzut pe Isus venind la el și a zis: “Iată Mielul lui Dumnezeu, care ridică păcatul lumii!

\VS{30}Acesta este Acela despre care am spus: “După mine vine un om care este preferat înaintea mea, pentru că a fost înaintea mea”.

\VS{31}Eu nu-l cunoșteam, dar pentru aceasta am venit botezând în apă, ca să fie descoperit lui Israel.”

\VS{32}Ioan a depus mărturie și a zis: “Am văzut Duhul Sfânt coborând din cer ca un porumbel și a rămas peste El.

\VS{33}Eu nu l-am recunoscut, dar cel care m-a trimis să botez în apă mi-a spus: “Pe oricine vei vedea Duhul coborând și rămânând asupra lui, acela este cel care botează în Duhul Sfânt.

\VS{34}Eu am văzut și am mărturisit că acesta este Fiul lui Dumnezeu.”


\par }{\PP \VS{35}A doua zi, Ioan, stând în picioare cu doi din ucenicii săi,

\VS{36}s-a uitat la Isus, care mergea, și a zis: “Iată Mielul lui Dumnezeu!”

\VS{37}Cei doi discipoli l-au auzit vorbind și l-au urmat pe Isus.

\VS{38}Isus s-a întors și, văzându-i că îi urmăreau, le-a zis:
{\WJ{“Ce căutați?”
}}

\par }{\PP I-au zis: “Rabi” (adică, în traducere, Învățător), “unde stai?”.


\par }{\PP \VS{39}Și le-a zis: “{\WJ{Veniți și vedeți.”
}}

\par }{\PP Au venit și au văzut unde stătea și au rămas cu el în acea zi. Era pe la ora zece.

\VS{40}Unul dintre cei doi care l-au auzit pe Ioan și l-au urmat a fost Andrei, fratele lui Simon Petru.

\VS{41}El l-a găsit mai întâi pe fratele său, Simon, și i-a spus: “L-am găsit pe Mesia!”. (ceea ce înseamnă, fiind interpretat, Hristos).

\VS{42}L-a adus la Isus. Isus s-a uitat la el și i-a spus: “{\WJ{Tu ești Simon, fiul lui Iona. Te vei numi Cefa” (ceea ce înseamnă, prin interpretare, Petru).
}}


\par }{\PP \VS{43}A doua zi, s-a hotărât să se ducă în Galileea și a găsit pe Filip. Isus i-a spus:
{\WJ{“Urmează-mă”.
}}

\VS{44}Filip era din Betsaida, cetatea lui Andrei și a lui Petru.

\VS{45}Filip l-a găsit pe Natanael și i-a spus: “L-am găsit pe Acela despre care au scris Moise în Lege, dar și profeții: Iisus din Nazaret, fiul lui Iosif”.


\par }{\PP \VS{46}Natanael I-a zis: “Poate ieși ceva bun din Nazaret?”

\par }{\PP Filip i-a zis: “Vino și vezi”.


\par }{\PP \VS{47}Isus a văzut pe Natanael venind la El și a zis despre el:
{\WJ{“Iată un israelit, în care nu este nici o înșelăciune!”
}}


\par }{\PP \VS{48}Natanael i-a zis: “De unde Mă cunoști?”

\par }{\PP Isus i-a răspuns: “{\WJ{Înainte ca Filip să te cheme, când erai sub smochin, te-am văzut.”
}}


\par }{\PP \VS{49}Natanael I-a răspuns: “Rabi, Tu ești Fiul lui Dumnezeu! Tu ești Regele lui Israel!”


\par }{\PP \VS{50}Isus i-a răspuns: “{\WJ{Pentru că ți-am spus: “Te-am văzut sub smochin”, crezi tu? Vei vedea lucruri mai mari decât acestea!”
}}

\VS{51}El i-a zis: “Cu
{\WJ{siguranță, vă spun tuturor că, după aceea, veți vedea cerul deschis și îngerii lui Dumnezeu urcând și coborând asupra Fiului Omului.”
}}



\par }\Chap{2}{\PP \VerseOne{1}A treia zi, a fost o nuntă în Cana Galileii. Mama lui Isus era acolo.

\VS{2}La nuntă a fost invitat și Isus, împreună cu discipolii săi.

\VS{3}Când s-a terminat vinul, mama lui Isus i-a spus: “Nu mai au vin”.


\par }{\PP \VS{4}Isus i-a zis: “{\WJ{Femeie, ce are aceasta de-a face cu tine și cu mine? Nu a venit încă ceasul Meu”.
}}


\par }{\PP \VS{5}Mama lui a zis slujitorilor: “Faceți tot ce vă va spune.”


\par }{\PP \VS{6}Și erau acolo șase vase de piatră pentru apă, așezate după obiceiul iudeilor, care conțineau câte doi sau trei metrii fiecare.

\VS{7}Isus le-a zis:
{\WJ{“Umpleți vasele cu apă”. Și le-au umplut până la refuz.
}}

\VS{8}El le-a zis:
{\WJ{“Acum scoateți puțin și duceți-l la conducătorul sărbătorii.” Și au luat-o.
}}

\VS{9}Când conducătorul ospățului a gustat apa devenită acum vin și nu știa de unde provine (dar servitorii care scoseseră apa știau), conducătorul ospățului l-a chemat pe mire

\VS{10}și i-a spus: “Toată lumea să servească mai întâi vinul bun, iar după ce oaspeții au băut din belșug, apoi pe cel mai rău. Tu ai păstrat vinul bun până acum!”.

\VS{11}Acest început al semnelor sale, Isus l-a făcut în Cana Galileii și și-a dezvăluit gloria, iar discipolii săi au crezut în el.


\par }{\PP \VS{12}După aceea, S-a coborât în Capernaum, împreună cu mama Sa, cu frații Săi și cu ucenicii Săi, și au rămas acolo câteva zile.


\par }{\PP \VS{13}Se apropia Paștele iudeilor, și Isus s-a suit la Ierusalim.

\VS{14}A găsit în templu pe cei care vindeau boi, oi și porumbei și pe schimbătorii de bani șezând.

\VS{15}El a făcut un bici din funii și i-a alungat pe toți din templu, atât pe oi, cât și pe boi; apoi a vărsat banii schimbătorilor și le-a răsturnat mesele.

\VS{16}Celor care vindeau porumbei le-a zis: “Scoateți
{\WJ{aceste lucruri de aici! Nu faceți din casa Tatălui meu o piață!”
}}

\VS{17}Ucenicii lui și-au amintit că era scris: “Zelul pentru casa voastră mă va mânca”.
\XREF{✡}{{\XO 2:17 }{\XT{Psalmul 69:9}}}


\par }{\PP \VS{18}Și Iudeii I-au răspuns: “Ce semn ne arăți, dacă faci aceste lucruri?”


\par }{\PP \VS{19}Isus le-a răspuns: “{\WJ{Distrugeți templul acesta și în trei zile îl voi ridica.”
}}


\par }{\PP \VS{20}Iudeii ziceau: “A fost nevoie de patruzeci și șase de ani ca să se zidească templul acesta! Îl veți ridica voi în trei zile?”

\VS{21}Dar el a vorbit despre templul trupului său.

\VS{22}De aceea, când a înviat din morți, discipolii lui și-au adus aminte că el spusese acest lucru și au crezut în Scriptură și în cuvântul pe care îl spusese Isus.


\par }{\PP \VS{23}Când a fost la Ierusalim, în timpul sărbătorii Paștelui, mulți au crezut în Numele Lui și au văzut semnele pe care le făcea.

\VS{24}Dar Isus nu li s-a încredințat lor, pentru că îi cunoștea pe toți

\VS{25}și pentru că nu avea nevoie ca cineva să dea mărturie despre om, pentru că El însuși știa ce este în om.



\par }\Chap{3}{\PP \VerseOne{1}Și era un om dintre farisei, numit Nicodim, un fruntaș al iudeilor.

\VS{2}El a venit la Isus noaptea și I-a zis: “Rabi, știm că ești un învățător venit de la Dumnezeu, căci nimeni nu poate face aceste semne pe care le faci Tu, dacă nu este Dumnezeu cu el.”


\par }{\PP \VS{3}Isus i-a răspuns: “{\WJ{Adevărat îți spun că, dacă cineva nu se naște din nou, nu
}}\FTNT{*}{{\FR 3:3 }{\FT{Cuvântul tradus “din nou” aici și în
{\CHAR{Ioan 3:7}} (ἄνωθεν) înseamnă, de asemenea, “din nou” și “de sus”.}}}{\WJ{poate vedea Împărăția lui Dumnezeu.”
}}


\par }{\PP \VS{4}Nicodim I-a zis: “Cum se poate naște un om când este bătrân? Poate să intre a doua oară în pântecele mamei sale și să se nască?”


\par }{\PP \VS{5}Isus a răspuns: “{\WJ{Adevărat vă spun că, dacă nu se naște cineva din apă și din Duh, nu poate intra în Împărăția lui Dumnezeu.
}}

\VS{6}{\WJ{Ceea ce este născut din carne este carne. Ceea ce este născut din Spirit este spirit.
}}

\VS{7}{\WJ{Nu vă mirați că v-am spus: “Trebuie să vă nașteți din nou”.
}}

\VS{8}{\WJ{Vântul}}\FTNT{†}{{\FR 3:8 }{\FT{Același cuvânt grecesc (πνεῦμα) înseamnă vânt, suflare și duh.}}}{\WJ{suflă unde vrea, iar voi îi auziți sunetul, dar nu știți de unde vine și încotro se îndreaptă. La fel este oricine este născut din Spirit.”
}}


\par }{\PP \VS{9}Nicodim I-a răspuns: “Cum pot fi acestea?”


\par }{\PP \VS{10}Isus i-a răspuns: “Ești
{\WJ{tu învățătorul lui Israel și nu înțelegi aceste lucruri?
}}

\VS{11}{\WJ{Adevărat îți spun că noi spunem ceea ce știm și dăm mărturie despre ceea ce am văzut, iar tu nu primești mărturia noastră.
}}

\VS{12}{\WJ{Dacă v-am spus lucruri pământești și nu credeți, cum veți crede dacă vă voi spune lucruri cerești?
}}

\VS{13}Nimeni nu s-a
{\WJ{înălțat la cer decât Cel care S-a coborât din cer, Fiul Omului, care este în ceruri.
}}

\VS{14}{\WJ{După cum Moise a înălțat șarpele în pustiu, tot așa trebuie să fie înălțat și Fiul Omului,
}}

\VS{15}pentru ca
{\WJ{oricine crede în el să nu piară, ci să aibă viață veșnică.
}}

\VS{16}{\WJ{Fiindcă atât de mult a iubit Dumnezeu lumea, încât a dat pe singurul Său Fiu născut,}}\FTNT{‡}{{\FR 3:16 }{\FT{Expresia “unicul născut” provine din cuvântul grecesc “μονογενη”, care este uneori tradus prin “unicul născut” sau “unicul și unicul”.}}}
{\WJ{pentru ca oricine crede în El să nu piară, ci să aibă viață veșnică.
}}

\VS{17}{\WJ{Căci Dumnezeu nu L-a trimis pe Fiul Său în lume ca să judece lumea, ci pentru ca lumea să fie mântuită prin El.
}}

\VS{18}{\WJ{Cine crede în El nu este judecat. Cel care nu crede a fost deja judecat, pentru că nu a crezut în numele singurului Fiu născut al lui Dumnezeu.
}}

\VS{19}{\WJ{Aceasta este judecata: că lumina a venit în lume și oamenii au iubit mai degrabă întunericul decât lumina, pentru că faptele lor erau rele.
}}

\VS{20}{\WJ{Căci oricine face răul urăște lumina și nu vine la lumină, ca să nu i se descopere faptele.
}}

\VS{21}{\WJ{Dar cel care face adevărul vine la lumină, pentru ca faptele lui să fie descoperite, ca să se vadă că au fost făcute în Dumnezeu.”
}}


\par }{\PP \VS{22}După aceste lucruri, Isus a venit cu ucenicii Săi în ținutul Iudeii. A rămas acolo cu ei și a botezat.

\VS{23}Ioan boteza și el în Enon, lângă Salim, pentru că acolo era multă apă. Ei au venit și au fost botezați;

\VS{24}căci Ioan nu fusese încă aruncat în închisoare.

\VS{25}De aceea, s-a iscat o dispută între discipolii lui Ioan și unii iudei în legătură cu purificarea.

\VS{26}Aceștia au venit la Ioan și i-au zis: “Rabi, cel care era cu tine dincolo de Iordan, despre care ai mărturisit, iată că el botează și toată lumea vine la el.”


\par }{\PP \VS{27}Ioan a răspuns: “Nimic nu poate primi omul, dacă nu-i este dat din cer.

\VS{28}Voi înșivă mărturisiți că am spus: “Nu eu sunt Hristosul”, ci “Eu am fost trimis înaintea lui”.

\VS{29}Cel care are mireasa este mirele; dar prietenul mirelui, care stă în picioare și îl aude, se bucură foarte mult din cauza glasului mirelui. De aceea bucuria mea este deplină.

\VS{30}El trebuie să crească, dar eu trebuie să mă micșorez.


\par }{\PP \VS{31}“Cel ce vine de sus este mai presus de toate. Cel care vine de pe pământ aparține pământului și vorbește despre pământ. Cel care vine din cer este mai presus de toate.

\VS{32}Ceea ce a văzut și a auzit, despre aceasta mărturisește; și nimeni nu primește mărturia lui.

\VS{33}Cel care a primit mărturia lui și-a pus pecetea pe aceasta: Dumnezeu este adevărat.

\VS{34}Căci cel pe care l-a trimis Dumnezeu vorbește cuvintele lui Dumnezeu; căci Dumnezeu dă Duhul fără măsură.

\VS{35}Tatăl îl iubește pe Fiul și a dat toate lucrurile în mâna lui.

\VS{36}Cine crede în Fiul are viață veșnică, dar cine nu\FTNT{§}{{\FR 3:36 }{\FT{Același cuvânt poate fi tradus în acest context prin “nu ascultă” sau “nu crede”.}}} ascultă de Fiul nu va vedea viața, ci mânia lui Dumnezeu rămâne asupra lui.”



\par }\Chap{4}{\PP \VerseOne{1}Când a aflat Domnul că fariseii au auzit că Isus face și botează mai mulți ucenici decât Ioan

\VS{2}(deși nu Isus însuși boteza, ci ucenicii lui),

\VS{3}a părăsit Iudeea și s-a dus în Galileea.

\VS{4}Trebuia să treacă prin Samaria.

\VS{5}Și a ajuns la o cetate din Samaria, numită Sihar, lângă parcela de pământ pe care Iacov o dăduse fiului său Iosif.

\VS{6}Acolo se afla fântâna lui Iacov. Isus, deci, fiind obosit de drum, s-a așezat lângă fântână. Era pe la ceasul al șaselea.
\FTNT{*}{{\FR 4:6 }{\FT{prânz}}}


\par }{\PP \VS{7}O femeie din Samaria a venit să scoată apă. Isus i-a zis:
{\WJ{“Dă-mi să beau”.
}}

\VS{8}Căci ucenicii Săi plecaseră în cetate să cumpere mâncare.


\par }{\PP \VS{9}Atunci femeia samariteancă i-a zis: “Cum se face că Tu, care ești iudeu, ceri de la mine, o samariteancă, să beau ceva?” (Căci iudeii nu au de-a face cu samaritenii).


\par }{\PP \VS{10}Isus i-a răspuns: “{\WJ{Dacă ai fi cunoscut darul lui Dumnezeu și cine este Cel ce-ți zice: “Dă-Mi să beau”, L-ai fi rugat și El ți-ar fi dat apă vie.”
}}


\par }{\PP \VS{11}Femeia i-a zis: “Domnule, nu ai cu ce să scoți apă, și fântâna este adâncă. Atunci de unde iei apa aceea vie?

\VS{12}Ești tu mai mare decât tatăl nostru Iacov, care ne-a dat fântâna și a băut el însuși din ea, la fel ca și copiii și vitele lui?”


\par }{\PP \VS{13}Isus i-a răspuns: “{\WJ{Oricine va bea din apa aceasta va înseta iarăși;
}}

\VS{14}{\WJ{dar oricine va bea din apa pe care i-o voi da Eu, nu va mai înseta niciodată, ci apa pe care i-o voi da Eu va deveni în el un izvor de apă care izvorăște pentru viața veșnică.”
}}


\par }{\PP \VS{15}Femeia I-a zis: “Dă-mi, Doamne, apa aceasta, ca să nu-mi fie sete și să nu vin până aici ca să beau.”


\par }{\PP \VS{16}Isus i-a zis:
{\WJ{“Du-te, cheamă pe bărbatul tău și vino aici.”
}}


\par }{\PP \VS{17}Femeia a răspuns: “Nu am bărbat.”

\par }{\PP Isus i-a zis: “{\WJ{Bine ai zis: “Nu am bărbat”,
}}

\VS{18}{\WJ{căci ai avut cinci soți; și cel pe care-l ai acum nu este bărbatul tău. Aceasta ai spus-o cu adevărat”.
}}


\par }{\PP \VS{19}Femeia i-a zis: “Domnule, văd că ești prooroc.

\VS{20}Părinții noștri s-au închinat în muntele acesta, iar voi, iudeii, spuneți că în Ierusalim este locul unde trebuie să se închine oamenii.”


\par }{\PP \VS{21}Isus i-a zis: “{\WJ{Femeie, crede-mă, vine ceasul când nici în muntele acesta, nici în Ierusalim nu vă veți închina Tatălui.
}}

\VS{22}{\WJ{Voi vă închinați la ceea ce nu cunoașteți. Noi ne închinăm la ceea ce cunoaștem, căci mântuirea vine de la iudei.
}}

\VS{23}{\WJ{Dar vine ceasul, și acum este, când adevărații închinători se vor închina Tatălui în duh și în adevăr, pentru că Tatăl caută astfel de închinători pentru a fi închinătorii Lui.
}}

\VS{24}{\WJ{Dumnezeu este duh, iar cei care i se închină trebuie să se închine în duh și în adevăr.”
}}


\par }{\PP \VS{25}Femeia I-a zis: “Știu că vine Mesia, Cel ce se numește Hristos.\FTNT{†}{{\FR 4:25 }{\FT{“Mesia” (ebraică) și “Hristos” (greacă) înseamnă amândouă “Unsul”.}}} Când va veni, ne va spune toate lucrurile.”


\par }{\PP \VS{26}Isus i-a zis: “{\WJ{Eu sunt Cel ce vorbește cu tine.”
}}


\par }{\PP \VS{27}În clipa aceea au venit ucenicii Lui. Ei s-au mirat că vorbea cu o femeie; dar nimeni nu i-a întrebat: “Ce cauți?” sau “De ce vorbești cu ea?”.

\VS{28}Femeia a lăsat vasul cu apă, s-a dus în cetate și a spus oamenilor:

\VS{29}“Veniți să vedeți un om care mi-a spus tot ce am făcut. Poate fi acesta Hristosul?”

\VS{30}Ei au ieșit din cetate și veneau la el.


\par }{\PP \VS{31}Între timp, ucenicii îl îndemnau, zicând: “Rabi, mănâncă.”


\par }{\PP \VS{32}Dar El le-a zis: “{\WJ{Am de mâncat o mâncare de care voi nu știți.”
}}


\par }{\PP \VS{33}Și ucenicii ziceau unul către altul: “I-a adus cineva ceva de mâncare?”


\par }{\PP \VS{34}Isus le-a zis: “{\WJ{Hrana Mea este să fac voia Celui ce M-a trimis și să împlinesc lucrarea Lui.
}}

\VS{35}{\WJ{Nu spuneți voi: “Mai sunt încă patru luni până la seceriș?”? Iată, vă spun: ridicați-vă ochii și priviți câmpurile, că ele sunt deja albe pentru seceriș.
}}

\VS{36}{\WJ{Cel care seceră primește plata și culege roade pentru viața veșnică, pentru ca atât cel care seamănă, cât și cel care seceră să se bucure împreună.
}}

\VS{37}{\WJ{Căci în aceasta este adevărată zicala: “Unul seamănă, iar altul culege”.
}}

\VS{38}{\WJ{Eu v-am trimis să secerați ceea ce nu ați muncit. Alții au muncit, iar tu ai intrat în munca lor.”
}}


\par }{\PP \VS{39}Mulți samariteni din cetatea aceea au crezut în El, din pricina cuvântului femeii care mărturisea: “El mi-a spus tot ce am făcut.”

\VS{40}Când au venit la el, samaritenii l-au rugat să rămână cu ei. El a rămas acolo două zile.

\VS{41}Mulți alții au crezut datorită cuvântului lui.

\VS{42}Ei i-au spus femeii: “Acum credem, nu din cauza vorbelor tale, ci pentru că am auzit cu ochii noștri și știm că acesta este cu adevărat Hristosul, Mântuitorul lumii.”


\par }{\PP \VS{43}După cele două zile, a ieșit de acolo și s-a dus în Galileea.

\VS{44}Căci Isus însuși a mărturisit că un profet nu are onoare în țara sa.

\VS{45}Când a ajuns în Galileea, galileenii L-au primit, după ce au văzut tot ce făcuse la Ierusalim la sărbătoare, căci și ei mergeau la sărbătoare.

\VS{46}Isus a venit deci din nou la Cana Galileii, unde a transformat apa în vin. Era un nobil al cărui fiu era bolnav la Capernaum.

\VS{47}Când a auzit că Isus a ieșit din Iudeea în Galileea, s-a dus la el și l-a rugat să se coboare și să-i vindece fiul, căci era la un pas de moarte.

\VS{48}Isus i-a zis:
{\WJ{“Dacă nu vei vedea semne și minuni, nu vei crede nicidecum”.
}}


\par }{\PP \VS{49}Și nobilul i-a zis: “Domnule, coboară-te, până nu moare copilul meu.”


\par }{\PP \VS{50}Isus i-a zis:
{\WJ{“Du-te și mergi. Fiul tău trăiește”. Omul a crezut cuvântul pe care i-l spusese Isus și a plecat.
}}

\VS{51}Pe când cobora, l-au întâlnit slujitorii lui și i-au raportat: “Fiul tău trăiește!”

\VS{52}Și i-a întrebat pe ei la ce oră a început să se facă bine. Ei i-au spus deci: “Ieri, în ceasul al șaptelea,\FTNT{‡}{{\FR 4:52 }{\FT{1:00 p.m.}}} febra l-a lăsat.”

\VS{53}Astfel, tatăl a știut că a fost în ceasul acela în care Isus i-a spus:
{\WJ{“Fiul tău trăiește.” El a crezut, ca și toată casa lui.
}}

\VS{54}Acesta este iarăși al doilea semn pe care l-a făcut Isus, după ce a ieșit din Iudeea în Galileea.



\par }\Chap{5}{\PP \VerseOne{1}După aceste lucruri, a fost o sărbătoare a iudeilor, și Isus s-a suit la Ierusalim.

\VS{2}În Ierusalim, lângă poarta oilor, se află un bazin care se numește în ebraică “Betesda”, având cinci pridvoare.

\VS{3}În acestea zăcea o mare mulțime de bolnavi, orbi, șchiopi sau paralizați, care așteptau să se miște apa;

\VS{4}pentru că un înger se cobora în anumite momente în bazin și agita apa. Cel care intra primul în bazin, după ce se agita apa, era vindecat de orice boală pe care o avea.
\FTNT{*}{{\FR 5:4 }{\FT{NU omite de la “așteptând” din versetul 3 până la sfârșitul versetului 4.}}}

\VS{5}Era acolo un om care era bolnav de treizeci și opt de ani.

\VS{6}Când Isus l-a văzut zăcând acolo și a știut că era bolnav de multă vreme, l-a întrebat:
{\WJ{“Vrei să te faci sănătos?”
}}


\par }{\PP \VS{7}Și bolnavul i-a răspuns: “Domnule, nu am pe nimeni care să mă bage în bazin când se agită apa, ci, în timp ce vin, altul coboară înaintea mea.”


\par }{\PP \VS{8}Isus i-a zis:
{\WJ{“Scoală-te, ia-ți rogojina și umblă.”
}}


\par }{\PP \VS{9}Îndată omul s-a făcut sănătos, și-a luat roata și a mers.

\par }{\PP În ziua aceea era Sabat.

\VS{10}Și iudeii au zis către cel vindecat: “Este Sabat. Nu-ți este îngăduit să cari covorul”.


\par }{\PP \VS{11}El le-a răspuns: “Cel ce m-a făcut sănătos mi-a zis: “{\WJ{Ia-ți roata și umblă!””
}}


\par }{\PP \VS{12}Și L-au întrebat: “Cine este omul care ți-a zis: “{\WJ{Ia-ți roata și umblă”?”
}}


\par }{\PP \VS{13}Dar cel vindecat nu știa cine era, pentru că Isus se retrăsese, fiindcă era mulțime în locul acela.


\par }{\PP \VS{14}După aceea, Isus l-a găsit în templu și i-a zis: “{\WJ{Iată că te-ai făcut sănătos. Nu mai păcătui, ca să nu ți se întâmple nimic mai rău.”
}}


\par }{\PP \VS{15}Omul s-a dus și a spus iudeilor că Isus îl făcuse sănătos.

\VS{16}Din această cauză, iudeii l-au prigonit pe Isus și căutau să-L omoare, pentru că făcea aceste lucruri în ziua de Sabat.

\VS{17}Dar Isus le-a răspuns:
{\WJ{“Tatăl Meu lucrează încă, așa că lucrez și eu.”
}}


\par }{\PP \VS{18}De aceea iudeii căutau cu atât mai mult să-L omoare, cu cât nu numai că a călcat Sabatul, dar a și numit pe Dumnezeu Tatăl său, făcându-se pe sine egal cu Dumnezeu.

\VS{19}Isus le-a răspuns: “{\WJ{Adevărat, vă spun că Fiul nu poate face nimic de la sine, ci doar ceea ce vede pe Tatăl făcând. Căci tot ceea ce face El, face și Fiul la fel.
}}

\VS{20}{\WJ{Căci Tatăl are afecțiune pentru Fiul și îi arată toate lucrurile pe care le face el însuși. Îi va arăta lucrări mai mari decât acestea, ca să vă minunați.
}}

\VS{21}{\WJ{Căci, după cum Tatăl înviază morții și le dă viață, tot așa și Fiul dă viață oricui dorește.
}}

\VS{22}{\WJ{Căci Tatăl nu judecă pe nimeni, ci a dat toată judecata Fiului,
}}

\VS{23}pentru ca
{\WJ{toți să cinstească pe Fiul, așa cum cinstesc pe Tatăl. Cine nu-l onorează pe Fiul, nu-l onorează pe Tatăl care l-a trimis.
}}


\par }{\PP \VS{24}“{\WJ{Adevărat vă spun că cine ascultă cuvântul Meu și crede pe Cel ce M-a trimis pe Mine are viață veșnică și nu vine la judecată, ci a trecut din moarte la viață.
}}

\VS{25}{\WJ{Adevărat vă spun că vine ceasul, și acum este, când cei morți vor auzi glasul Fiului lui Dumnezeu, iar cei ce-l vor auzi vor trăi.
}}

\VS{26}{\WJ{Căci, după cum Tatăl are viața în Sine, tot așa a dat și Fiului să aibă viața în Sine.
}}

\VS{27}{\WJ{I-a dat, de asemenea, autoritatea de a face judecată, pentru că este fiu al omului.
}}

\VS{28}{\WJ{Nu vă mirați de aceasta, căci vine ceasul în care toți cei care sunt în morminte vor auzi glasul lui
}}

\VS{29}{\WJ{și vor ieși: cei care au făcut binele, la învierea vieții, iar cei care au făcut răul, la învierea judecății.
}}

\VS{30}{\WJ{Eu nu pot de la mine însumi să fac nimic. Cum aud, judec; și judecata mea este dreaptă, pentru că nu caut voia mea, ci voia Tatălui meu care m-a trimis.
}}


\par }{\PP \VS{31}{\WJ{Dacă mărturisesc despre mine însumi, mărturia mea nu este valabilă.
}}

\VS{32}Cel
{\WJ{care depune mărturie despre mine este altul. Știu că mărturia pe care o depune el despre mine este adevărată.
}}

\VS{33}{\WJ{Tu ai trimis la Ioan și el a mărturisit adevărul.
}}

\VS{34}{\WJ{Dar mărturia pe care o primesc eu nu vine de la om. Totuși, spun aceste lucruri pentru ca voi să fiți mântuiți.
}}

\VS{35}{\WJ{El era lampa aprinsă și strălucitoare și voi ați vrut să vă bucurați pentru o vreme de lumina lui.
}}

\VS{36}{\WJ{Dar mărturia pe care o am eu este mai mare decât cea a lui Ioan; pentru că lucrările pe care Tatăl mi-a dat să le îndeplinesc, chiar lucrările pe care le fac eu, mărturisesc despre mine că Tatăl m-a trimis.
}}

\VS{37}{\WJ{Însuși Tatăl, care m-a trimis, a mărturisit despre mine. Voi nu i-ați auzit niciodată vocea și nici nu i-ați văzut forma.
}}

\VS{38}{\WJ{Nu aveți cuvântul lui care trăiește în voi, pentru că nu credeți în cel pe care l-a trimis.
}}


\par }{\PP \VS{39}{\WJ{Voi cercetați Scripturile, pentru că credeți că în ele aveți viața veșnică; și acestea sunt cele ce mărturisesc despre Mine.
}}

\VS{40}{\WJ{Și totuși nu vreți să veniți la Mine, ca să aveți viața.
}}

\VS{41}{\WJ{Eu nu primesc slavă de la oameni.
}}

\VS{42}{\WJ{Dar vă cunosc pe voi, că nu aveți în voi înșivă dragostea lui Dumnezeu.
}}

\VS{43}{\WJ{Eu am venit în numele Tatălui Meu, și voi nu Mă primiți. Dacă vine altul în numele lui, îl veți primi.
}}

\VS{44}{\WJ{Cum puteți crede voi, care primiți slavă unii de la alții, iar voi nu căutați slava care vine de la singurul Dumnezeu?
}}


\par }{\PP \VS{45}“Să
{\WJ{nu credeți că vă voi acuza la Tatăl. Este unul care vă acuză, chiar Moise, în care v-ați pus nădejdea.
}}

\VS{46}{\WJ{Căci, dacă l-ați crede pe Moise, m-ați crede și pe mine, căci el a scris despre mine.
}}

\VS{47}{\WJ{Dardacă nu credeți scrierile lui, cum veți crede cuvintele mele?”.
}}



\par }\Chap{6}{\PP \VerseOne{1}După acestea, Isus a plecat de cealaltă parte a Mării Galileii, care se numește Marea Tiberiadei.

\VS{2}O mare mulțime de oameni Îl urmau, pentru că vedeau semnele pe care le făcea asupra celor bolnavi.

\VS{3}Isus s-a urcat pe munte și a șezut acolo cu discipolii Săi.

\VS{4}Se apropia Paștele, sărbătoarea iudeilor.

\VS{5}Isus, ridicându-și ochii și văzând că o mare mulțime venea la El, a zis lui Filip:
{\WJ{“De unde vom cumpăra pâine, ca să mănânce aceștia?”
}}

\VS{6}A spus aceasta ca să-l pună la încercare, căci știa el însuși ce va face.


\par }{\PP \VS{7}Filip i-a răspuns: “Nu le ajunge pâinea de două sute de denari\FTNT{*}{{\FR 6:7 }{\FT{Un denar era o monedă de argint care valora aproximativ o zi de salariu pentru un muncitor agricol, astfel că 200 de denari ar reprezenta între 6 și 7 luni de salariu.}}}, ca să primească fiecare din ei câte puțin.”


\par }{\PP \VS{8}Unul din ucenicii Săi, Andrei, fratele lui Simon Petru, i-a zis:

\VS{9}“Este aici un băiat care are cinci pâini de orz și doi pești; dar ce sunt aceștia între atâția?”


\par }{\PP \VS{10}Isus a zis:
{\WJ{“Fă să șadă poporul.” Or, în locul acela era multă iarbă. Deci oamenii au luat loc, în număr de aproximativ cinci mii.
}}

\VS{11}Isus a luat pâinile și, după ce a mulțumit, a împărțit ucenicilor, iar ucenicii celor ce stăteau jos, la fel și din pește, atât cât au dorit.

\VS{12}După ce s-au săturat, a zis ucenicilor Săi: “Adunați
{\WJ{bucățile rămase, ca să nu se piardă nimic.”
}}

\VS{13}Ei le-au adunat și au umplut douăsprezece coșuri cu bucățile rupte din cele cinci pâini de orz, care rămăseseră de la cei care mâncaseră.

\VS{14}De aceea, când a văzut poporul semnul pe care-l făcea Isus, a zis: “Acesta este cu adevărat profetul care vine în lume.”

\VS{15}Isus, deci, dându-și seama că sunt pe cale să vină și să îl ia cu forța ca să îl facă rege, s-a retras din nou pe munte, de unul singur.


\par }{\PP \VS{16}Când s-a făcut seară, ucenicii Lui s-au coborât la mare.

\VS{17}S-au urcat în corabie și se îndreptau peste mare spre Capernaum. Era deja întuneric și Isus nu venise la ei.

\VS{18}Marea era agitată de un vânt mare care sufla.

\VS{19}După ce au vâslit, așadar, vreo douăzeci și cinci sau treizeci de stadii,\FTNT{†}{{\FR 6:19 }{\FT{25 și 30 de stadii înseamnă aproximativ 5 până la 6 kilometri sau aproximativ 3 până la 4 mile.}}} au văzut pe Isus mergând pe mare\FTNT{‡}{{\FR 6:19 }{\FT{Vezi Iov 9:8}}} și apropiindu-se de barcă; și s-au temut.

\VS{20}Dar El le-a zis:
{\WJ{“Eu sunt;}}\FTNT{§}{{\FR 6:20 }{\FT{sau, EU SUNT}}}
{\WJ{nu vă temeți.”
}}

\VS{21}Ei au fost deci dispuși să îl primească în barcă. Imediat, barca a ajuns la țărmul spre care se îndreptau.


\par }{\PP \VS{22}A doua zi, mulțimea care stătea de cealaltă parte a mării a văzut că nu mai era acolo altă corabie decât cea în care se îmbarcaseră ucenicii Lui și că Isus nu intrase cu ucenicii Lui în corabie, ci ucenicii Lui plecaseră singuri.

\VS{23}Cu toate acestea, niște bărci venite dinspre Tiberiada s-au apropiat de locul unde mâncaseră pâinea după ce Domnul a mulțumit.

\VS{24}Așadar, când mulțimea a văzut că Isus și discipolii lui nu erau acolo, s-a urcat ea însăși în bărci și a venit la Capernaum, căutându-l pe Isus.

\VS{25}Când L-au găsit pe partea cealaltă a mării, L-au întrebat: “Rabi, când ai venit aici?”


\par }{\PP \VS{26}Isus le-a răspuns: “{\WJ{Adevărat vă spun că Mă căutați, nu pentru că ați văzut semne, ci pentru că ați mâncat din pâini și v-ați săturat.
}}

\VS{27}{\WJ{Nu lucrați pentru hrana care piere, ci pentru hrana care rămâne pentru viața veșnică, pe care v-o va da Fiul Omului. Căci Dumnezeu Tatăl l-a pecetluit.”
}}


\par }{\PP \VS{28}Și I-au zis: “Ce trebuie să facem ca să lucrăm lucrările lui Dumnezeu?”


\par }{\PP \VS{29}Isus le-a răspuns: “{\WJ{Aceasta este lucrarea lui Dumnezeu: să credeți în Cel pe care L-a trimis El.”
}}


\par }{\PP \VS{30}Și I-au zis: “Ce faci Tu ca semn, ca să te vedem și să te credem? Ce lucrare faci tu?

\VS{31}Părinții noștri au mâncat mană în pustiu. După cum este scris: “Le-a dat să mănânce pâine din cer”.”
\FTNT{*}{{\FR 6:31 }{\FT{greacă și în ebraică se folosește același cuvânt pentru “cer”, “ceruri”, “cer” și “aer”.}}}\FTNT{†}{{\FR 6:31 }{\FT{Exodul 16:4; Neemia 9:15;
{\CHAR{Psalmul 78:24-25}}}}}


\par }{\PP \VS{32}Isus le-a zis: “{\WJ{Adevărat vă spun că nu Moise v-a dat pâinea din cer, ci Tatăl Meu vă dă adevărata pâine din cer.
}}

\VS{33}{\WJ{Căci pâinea lui Dumnezeu este cea care se coboară din cer și dă viață lumii.”
}}


\par }{\PP \VS{34}Și I-au zis: “Doamne, dă-ne totdeauna pâinea aceasta.”


\par }{\PP \VS{35}Isus le-a zis: “{\WJ{Eu sunt pâinea vieții. Cine vine la mine nu va flămânzi și cine crede în mine nu va înseta niciodată.
}}

\VS{36}{\WJ{Dar v-am spus că m-ați văzut și totuși nu credeți.
}}

\VS{37}{\WJ{Toți cei pe care mi-i dă Tatăl vor veni la mine. Pe cel care vine la mine nu-l voi da afară în niciun fel.
}}

\VS{38}{\WJ{Căci Eu m-am coborât din cer, nu ca să fac voia Mea, ci voia Celui care M-a trimis.
}}

\VS{39}{\WJ{Aceasta este voia Tatălui Meu care m-a trimis: ca din toți cei pe care Mi i-a dat să nu pierd nimic, ci să-i înviez în ziua de apoi.
}}

\VS{40}{\WJ{Aceasta este voia Celui care m-a trimis: ca oricine vede pe Fiul și crede în El să aibă viața veșnică; și Eu îl voi învia în ziua de apoi.”
}}


\par }{\PP \VS{41}De aceea Iudeii cârteau împotriva Lui, pentru că zicea:
{\WJ{“Eu sunt pâinea care S-a pogorât din cer”.
}}

\VS{42}Ei ziceau: “Nu este acesta Isus, fiul lui Iosif, al cărui tată și mamă îi știm? Cum de spune atunci:
{\WJ{“Eu am coborât din cer?”?”
}}


\par }{\PP \VS{43}De aceea Isus le-a răspuns: “{\WJ{Nu vă certați între voi.
}}

\VS{44}Nimeni
{\WJ{nu poate veni la Mine, dacă nu-l atrage Tatăl care M-a trimis; și Eu îl voi învia în ziua de apoi.
}}

\VS{45}{\WJ{Este scris în profeți: “Toți vor fi învățați de Dumnezeu”.
}}\FTNT{‡}{{\FR 6:45 }{\FT{{\CHAR{Isaia 54:13}}}}}{\WJ{De aceea, oricine aude de la Tatăl și a învățat, vine la mine.
}}

\VS{46}{\WJ{Nu că cineva a văzut pe Tatăl, decât cel care vine de la Dumnezeu. El l-a văzut pe Tatăl.
}}

\VS{47}{\WJ{Adevărat, vă spun, cel ce crede în Mine are viața veșnică.
}}

\VS{48}{\WJ{Eu sunt pâinea vieții.
}}

\VS{49}{\WJ{Părinții voștri au mâncat mană în pustiu și au murit.
}}

\VS{50}{\WJ{Aceasta este pâinea care se coboară din cer, ca oricine să mănânce din ea și să nu moară.
}}

\VS{51}{\WJ{Eu sunt pâinea vie care s-a coborât din cer. Dacă cineva mănâncă din pâinea aceasta, va trăi în veci. Da, pâinea pe care o voi da pentru viața lumii este carnea mea.”
}}


\par }{\PP \VS{52}Și iudeii se certau între ei, zicând: “Cum ne va da Acesta să mâncăm carnea Lui?”


\par }{\PP \VS{53}Isus le-a zis: “{\WJ{Adevărat vă spun că, dacă nu mâncați carnea Fiului Omului și nu beți sângele Lui, nu aveți viață în voi înșivă.
}}

\VS{54}{\WJ{Cine mănâncă carnea Mea și bea sângele Meu are viață veșnică și Eu îl voi învia în ziua de apoi.
}}

\VS{55}{\WJ{Căci carnea Mea este cu adevărat hrană și sângele Meu este cu adevărat băutură.
}}

\VS{56}{\WJ{Cine mănâncă carnea Mea și bea sângele Meu trăiește în Mine și Eu în el.
}}

\VS{57}{\WJ{După cum Tatăl cel viu m-a trimis pe mine și eu trăiesc datorită Tatălui, tot așa și cel care se hrănește din mine va trăi datorită mie.
}}

\VS{58}{\WJ{Aceasta este pâinea care s-a coborât din cer — nu cum au mâncat părinții noștri mana și au murit. Cel care mănâncă această pâine va trăi în veci.”
}}

\VS{59}El spunea aceste lucruri în sinagogă, pe când învăța în Capernaum.


\par }{\PP \VS{60}De aceea mulți dintre ucenicii Lui, auzind aceasta, au zis: “Greu este să spui aceasta! Cine poate să o asculte?”


\par }{\PP \VS{61}Dar Isus, știind că ucenicii Săi murmurau din pricina aceasta, le-a zis: “{\WJ{Oare vă împiedică aceasta?
}}

\VS{62}{\WJ{Atunci, dacă ați vrea să vedeți pe Fiul Omului înălțându-se la locul unde era mai înainte?
}}

\VS{63}{\WJ{Spiritul este cel care dă viață. Carnea nu folosește la nimic. Cuvintele pe care vi le spun sunt spirit și sunt viață.
}}

\VS{64}{\WJ{Dar sunt unii dintre voi care nu cred.”}} Pentru că Isus știa de la început cine sunt cei care nu credeau și cine sunt cei care aveau să-l trădeze.

\VS{65}El a spus: “De
{\WJ{aceea v-am spus că nimeni nu poate veni la mine, dacă nu-i este dat de Tatăl meu.”
}}


\par }{\PP \VS{66}Mulți din ucenicii Lui s-au întors și nu mai umblau cu El.

\VS{67}Isus le-a zis deci celor doisprezece: “{\WJ{Nu cumva vreți și voi să plecați?”
}}


\par }{\PP \VS{68}Simon Petru I-a răspuns: “Doamne, la cine să mergem? Tu ai cuvintele vieții veșnice.

\VS{69}Noi am ajuns să credem și să știm că Tu ești Hristosul, Fiul Dumnezeului celui viu.”


\par }{\PP \VS{70}Isus le-a răspuns: “{\WJ{Nu v-am ales Eu pe voi, cei doisprezece, și unul dintre voi este un diavol?”
}}

\VS{71}Și vorbea despre Iuda, fiul lui Simon Iscarioteanul, căci el era cel care avea să-L trădeze, fiind unul dintre cei doisprezece.



\par }\Chap{7}{\PP \VerseOne{1}După acestea, Isus umbla prin Galileea, căci nu voia să umble în Iudeea, pentru că iudeii căutau să-L omoare.

\VS{2}Or, sărbătoarea iudeilor, sărbătoarea Bobotezei, era aproape.

\VS{3}Frații Lui i-au zis: “Pleacă de aici și du-te în Iudeea, ca să vadă și ucenicii Tăi lucrările pe care le faci.

\VS{4}Căci nimeni nu face nimic în secret, în timp ce caută să fie cunoscut în mod deschis. Dacă tu faci aceste lucruri, descoperă-te lumii.”

\VS{5}Căci nici măcar frații lui nu credeau în el.


\par }{\PP \VS{6}Atunci Isus le-a zis: “{\WJ{Încă nu a venit vremea Mea, dar vremea voastră este totdeauna gata.
}}

\VS{7}{\WJ{Lumea nu vă poate urî pe voi, dar pe Mine Mă urăște, pentru că Eu mărturisesc despre ea că faptele ei sunt rele.
}}

\VS{8}{\WJ{Voi vă urcați la ospăț. Eu nu mă sui încă la acest ospăț, pentru că timpul meu nu s-a împlinit încă.”
}}


\par }{\PP \VS{9}După ce le-a spus aceste lucruri, a rămas în Galileea.

\VS{10}Când frații lui s-au suit la ospăț, s-a suit și el, dar nu în public, ci ca în secret.

\VS{11}Iudeii îl căutau deci la ospăț și întrebau: “Unde este?”

\VS{12}Și a fost multă murmurare în mulțime cu privire la el. Unii spuneau: “Este un om bun”. Alții spuneau: “Nu-i așa, dar el duce mulțimea în rătăcire”.

\VS{13}Totuși, nimeni nu vorbea deschis despre el, de frica iudeilor.

\VS{14}Dar, când era deja mijlocul sărbătorii, Isus s-a urcat în templu și învăța.

\VS{15}Iudeii se mirau deci și ziceau: “Cum de știe acesta litere, fără să fi fost educat niciodată?”


\par }{\PP \VS{16}Isus le-a răspuns: “{\WJ{Nu a Mea este învățătura Mea, ci a Celui ce M-a trimis.
}}

\VS{17}{\WJ{Dacă cineva vrea să facă voia Lui, va cunoaște învățătura, dacă este de la Dumnezeu sau dacă vorbesc de la mine însumi.
}}

\VS{18}{\WJ{Cine vorbește de la sine caută propria slavă, dar cine caută slava celui care l-a trimis pe el este adevărat și nu există nedreptate în el.
}}

\VS{19}{\WJ{Nu v-a dat Moise Legea, și totuși niciunul dintre voi nu respectă Legea? De ce căutați să mă omorâți?”
}}


\par }{\PP \VS{20}Mulțimea a răspuns: “Ai un demon! Cine caută să te ucidă?”


\par }{\PP \VS{21}Isus le-a răspuns: “{\WJ{Eu am făcut o singură lucrare și voi toți vă mirați de ea.
}}

\VS{22}{\WJ{Moise v-a dat circumcizia (nu că ar fi de la Moise, ci de la părinți) și în Sabat circumcideți un băiat.
}}

\VS{23}{\WJ{Dacă un băiat primește circumcizia în Sabat, ca să nu se încalce Legea lui Moise, vă supărați pe Mine pentru că am făcut un om complet sănătos în Sabat?
}}

\VS{24}{\WJ{Nu judecați după aparențe, ci judecați cu dreaptă judecată.”
}}


\par }{\PP \VS{25}De aceea, unii din cei din Ierusalim au zis: “Nu cumva nu este acesta pe care caută să-L omoare?

\VS{26}Iată că el vorbește deschis, și ei nu-i spun nimic. Se poate oare ca într-adevăr conducătorii să știe că acesta este cu adevărat Hristosul?

\VS{27}Totuși, noi știm de unde vine omul acesta, dar când va veni Hristosul, nimeni nu va ști de unde vine.”


\par }{\PP \VS{28}Isus a strigat în templu, învățând și zicând: “{\WJ{Voi Mă cunoașteți și știți de unde sunt. Eu n-am venit de la mine însumi, ci Cel care m-a trimis este adevărat, pe care voi nu-l cunoașteți.
}}

\VS{29}{\WJ{Îl cunosc, pentru că Eu sunt de la El și El m-a trimis.”
}}


\par }{\PP \VS{30}Au căutat deci să-L prindă, dar nimeni n-a pus mâna pe El, pentru că nu venise încă ceasul Lui.

\VS{31}Dar, din mulțime, mulți au crezut în El. Ei spuneau: “Când va veni Hristosul, nu va face mai multe semne decât cele pe care le-a făcut acesta, nu-i așa?”

\VS{32}Fariseii au auzit mulțimea murmurând aceste lucruri despre el și preoții cei mai de seamă și fariseii au trimis ofițeri ca să-l aresteze.


\par }{\PP \VS{33}Atunci Isus a zis: “{\WJ{Mai stau puțin cu voi, apoi Mă duc la Cel ce M-a trimis.
}}

\VS{34}{\WJ{Mă veți căuta și nu mă veți găsi. Nu puteți veni acolo unde sunt eu.”
}}


\par }{\PP \VS{35}Și iudeii ziceau între ei: “Unde se va duce omul acesta, ca să nu-l găsim? Se va duce oare în Dispersia dintre greci și îi va învăța pe greci?

\VS{36}Ce este cuvântul acesta pe care l-a spus:
{\WJ{“Mă veți căuta și nu Mă veți găsi”}} și
{\WJ{“Unde sunt Eu, nu puteți veni”}}?”


\par }{\PP \VS{37}În ultima și cea mai mare zi a praznicului, Isus, stând în picioare, a strigat: “{\WJ{Dacă însetează cineva, să vină la Mine și să bea!
}}

\VS{38}{\WJ{Cine crede în mine, după cum spune Scriptura, dinlăuntrul lui vor curge râuri de apă vie.”
}}

\VS{39}Dar a spus următoarele despre Duhul Sfânt, pe care urmau să-l primească cei care credeau în el. Căci Duhul Sfânt nu fusese încă dat, pentru că Isus nu fusese încă glorificat.


\par }{\PP \VS{40}Mulți din mulțime, auzind aceste cuvinte, ziceau: “Acesta este cu adevărat proorocul”.

\VS{41}Alții ziceau: “Acesta este Hristosul”. Dar unii ziceau: “Ce, oare Hristosul vine din Galileea?

\VS{42}Nu a spus oare Scriptura că Hristosul vine din neamul\FTNT{*}{{\FR 7:42 }{\FT{sau, sămânță}}} lui David și\FTNT{†}{{\FR 7:42 }{\FT{7:42 2 Samuel 7:12}}} din Betleem,\FTNT{‡}{{\FR 7:42 }{\FT{Mica 5:2}}} satul în care era David?”

\VS{43}Și s-a creat o dezbinare în mulțime din cauza lui.

\VS{44}Unii dintre ei ar fi vrut să-l aresteze, dar nimeni nu a pus mâna pe el.

\VS{45}Ofițerii au venit deci la preoții cei mai de seamă și la farisei, care i-au întrebat: “De ce nu l-ați adus?”


\par }{\PP \VS{46}Ofițerii au răspuns: “Nimeni n-a vorbit vreodată ca omul acesta!”


\par }{\PP \VS{47}Fariseii le-au răspuns: “Nu cumva și voi sunteți rătăciți?

\VS{48}A crezut oare vreunul dintre dregători sau vreunul dintre farisei în el?

\VS{49}Dar mulțimea aceasta care nu cunoaște Legea este blestemată.”


\par }{\PP \VS{50}Nicodim, care venise noaptea la El, fiind unul dintre ei, le-a zis:

\VS{51}“Oare legea noastră judecă pe cineva dacă nu-l aude mai întâi și nu știe ce face?”


\par }{\PP \VS{52}Ei I-au răspuns: “Și Tu ești din Galileea? Caută și vezi că din Galileea nu s-a ridicat niciun profet.”
\FTNT{§}{{\FR 7:52 }{\FT{Vezi
{\CHAR{Isaia 9:1}};
{\CHAR{Matei 4:13-16}}}}}


\par }{\PP \VS{53}Fiecare s-a dus la casa lui,



\par }\Chap{8}{\PP \VerseOne{1}Dar Isus s-a dus pe Muntele Măslinilor.


\par }{\PP \VS{2}Și, dis-de-dimineață, a intrat din nou în Templu și tot poporul a venit la El. El a șezut jos și i-a învățat.

\VS{3}Cărturarii și fariseii au adus o femeie prinsă în adulter. După ce au așezat-o la mijloc,

\VS{4}i-au spus: “Învățătorule, am găsit-o pe această femeie în adulter, chiar în flagrant.

\VS{5}Or, în legea noastră, Moise ne-a poruncit să ucidem cu pietre astfel de femei.\FTNT{*}{{\FR 8:5 }{\FT{Leviticul 20:10; Deuteronomul 22:22}}} Ce spui, așadar, despre ea?”

\VS{6}Ei au spus aceasta punându-l la încercare, ca să aibă de ce să-l acuze.

\par }{\PP Dar Isus s-a aplecat și a scris cu degetul pe pământ.

\VS{7}Dar, când ei continuau să-L întrebe, El și-a ridicat privirea și le-a zis: “{\WJ{Cine dintre voi este fără de păcat, să arunce primul cu piatra în ea.”
}}

\VS{8}Și iarăși s-a aplecat și a scris cu degetul pe pământ.


\par }{\PP \VS{9}Când au auzit, ei, încredințați de conștiința lor, au ieșit unul câte unul, începând de la cel mai bătrân până la cel din urmă. Isus a rămas singur cu femeia acolo unde era, în mijloc.

\VS{10}Isus, ridicându-se în picioare, a văzut-o și a zis:
{\WJ{“Femeie, unde sunt acuzatorii tăi? Nu te-a condamnat nimeni?”.
}}


\par }{\PP \VS{11}Ea a răspuns: “Nimeni, Doamne.”

\par }{\PP Isus a spus: “{\WJ{Nici Eu nu vă condamn. Mergeți pe drumul vostru. De acum încolo, să nu mai păcătuiești”.
}}\FTNT{†}{{\FR 8:11 }{\FT{NU include Ioan 7:53-Ioan 8:11, dar pune paranteze în jurul ei pentru a indica faptul că criticii textuali aveau mai puțină încredere în faptul că acesta era original.}}}


\par }{\PP \VS{12}Isus le-a vorbit din nou, zicând: “{\WJ{Eu sunt lumina lumii.}}\XREF{✡}{{\XO 8:12 }{\XT{Isaia 60:1}}}
{\WJ{Cine Mă urmează pe Mine nu va umbla în întuneric, ci va avea lumina vieții.”
}}


\par }{\PP \VS{13}Fariseii I-au zis: “Tu mărturisești despre tine însuți. Mărturia ta nu este valabilă”.


\par }{\PP \VS{14}Isus le-a răspuns: “{\WJ{Chiar dacă mărturisesc despre mine însumi, mărturia mea este adevărată, căci eu știu de unde am venit și unde mă duc; dar voi nu știți de unde am venit și unde mă duc.
}}

\VS{15}{\WJ{Voi judecați după trup. Eu nu judec pe nimeni.
}}

\VS{16}{\WJ{Chiar dacă judec, judecata mea este adevărată, pentru că nu sunt singur, ci sunt cu Tatăl care m-a trimis.
}}

\VS{17}De
{\WJ{asemenea, este scris în legea voastră că mărturia a doi oameni este valabilă.
}}\FTNT{‡}{{\FR 8:17 }{\FT{Deuteronomul 17:6; 19:15}}}

\VS{18}{\WJ{Eu sunt cel care mărturisește despre mine însumi, iar Tatăl care m-a trimis mărturisește despre mine.”
}}


\par }{\PP \VS{19}Și I-au zis: “Unde este Tatăl Tău?”

\par }{\PP Isus a răspuns: “{\WJ{Nu Mă cunoașteți nici pe Mine, nici pe Tatăl Meu. Dacă M-ați cunoaște pe Mine, L-ați cunoaște și pe Tatăl Meu.”
}}

\VS{20}Isus a rostit aceste cuvinte în vistierie, în timp ce învăța în templu. Dar nimeni nu l-a arestat, pentru că nu venise încă ceasul lui.

\VS{21}Isus le-a spus deci din nou: “Mă
{\WJ{duc, și voi mă veți căuta și veți muri în păcatele voastre. Unde Mă duc Eu, voi nu puteți veni”.
}}


\par }{\PP \VS{22}Iudeii ziceau: “Oare se va omorî El însuși, pentru că zice: “{\WJ{Unde Mă duc Eu, voi nu puteți veni”?”
}}


\par }{\PP \VS{23}El le-a zis: “{\WJ{Voi sunteți de jos. Eu sunt de sus. Voi sunteți din lumea aceasta. Eu nu sunt din lumea aceasta.
}}

\VS{24}{\WJ{De aceea v-am spus că veți muri în păcatele voastre; căci, dacă nu credeți că Eu sunt, veți muri în păcatele voastre.”
}}


\par }{\PP \VS{25}Și I-au zis: “Cine ești Tu?”

\par }{\PP Isus le-a spus: “{\WJ{Exact ceea ce v-am spus de la început.
}}

\VS{26}{\WJ{Am multe lucruri de spus și de judecat cu privire la voi. Totuși, cel care m-a trimis este adevărat; și lucrurile pe care le-am auzit de la el, acestea le spun lumii.”
}}


\par }{\PP \VS{27}Ei nu înțelegeau că El le vorbea despre Tatăl.

\VS{28}Isus le-a zis deci: “{\WJ{Când veți ridica pe Fiul Omului, atunci veți cunoaște că Eu sunt El și nu fac nimic de la mine însumi, ci spun aceste lucruri așa cum M-a învățat Tatăl Meu.
}}

\VS{29}{\WJ{Cel care m-a trimis este cu mine. Tatăl nu m-a lăsat singur, pentru că eu fac întotdeauna lucrurile care îi sunt plăcute.”
}}


\par }{\PP \VS{30}Pe când vorbea El aceste lucruri, mulți au crezut în El.

\VS{31}Isus a zis deci iudeilor care crezuseră în El: “{\WJ{Dacă rămâneți în cuvântul Meu, sunteți cu adevărat ucenicii Mei.
}}

\VS{32}{\WJ{Veți cunoaște adevărul și adevărul vă va face liberi.”
}}


\par }{\PP \VS{33}Ei i-au răspuns: “Noi suntem urmașii lui Avraam și n-am fost niciodată robi nimănui. Cum spui tu:
{\WJ{“Veți fi liberi”?”
}}


\par }{\PP \VS{34}Isus le-a răspuns: “{\WJ{Adevărat vă spun că oricine săvârșește un păcat este robul păcatului.
}}

\VS{35}{\WJ{O roabă nu locuiește în casă pentru totdeauna. Un fiu rămâne pentru totdeauna.
}}

\VS{36}{\WJ{Așadar, dacă Fiul vă face liberi, veți fi cu adevărat liberi.
}}

\VS{37}{\WJ{Știu că sunteți urmașii lui Avraam, dar voi căutați să mă omorâți, pentru că cuvântul meu nu-și găsește locul în voi.
}}

\VS{38}{\WJ{Eu spun lucrurile pe care le-am văzut la Tatăl meu, iar voi faceți și voi lucrurile pe care le-ați văzut la tatăl vostru.”
}}


\par }{\PP \VS{39}Ei i-au răspuns: “Tatăl nostru este Avraam.”

\par }{\PP Isus le-a zis: “{\WJ{Dacă ați fi copiii lui Avraam, ați face faptele lui Avraam.
}}

\VS{40}{\WJ{Dar acum căutați să Mă omorâți pe Mine, un om care v-a spus adevărul pe care l-am auzit de la Dumnezeu. Avraam nu a făcut asta.
}}

\VS{41}{\WJ{Voi faceți faptele tatălui vostru.”
}}

\par }{\PP Ei i-au spus: “Noi nu ne-am născut din imoralitate sexuală. Noi avem un singur Tată, Dumnezeu”.


\par }{\PP \VS{42}Isus le-a zis: “{\WJ{Dacă Dumnezeu ar fi tatăl vostru, M-ați iubi, căci Eu am ieșit și am venit de la Dumnezeu. Căci nu am venit de la mine însumi, ci el m-a trimis.
}}

\VS{43}{\WJ{De ce nu înțelegeți vorbirea mea? Pentru că nu puteți auzi cuvântul meu.
}}

\VS{44}{\WJ{Voi sunteți de la tatăl vostru, diavolul, și vreți să împliniți dorințele tatălui vostru. El a fost ucigaș de la început și nu stă în adevăr, pentru că nu este adevăr în el. Când vorbește o minciună, vorbește de unul singur, căci este un mincinos și tatăl minciunii.
}}

\VS{45}{\WJ{Dar pentru că spun adevărul, voi nu mă credeți.
}}

\VS{46}{\WJ{Care dintre voi mă convinge de păcat? Dacă spun adevărul, de ce nu mă credeți?
}}

\VS{47}{\WJ{Cine este de la Dumnezeu ascultă cuvintele lui Dumnezeu. De aceea voi nu auziți, pentru că nu sunteți din Dumnezeu.”
}}


\par }{\PP \VS{48}Atunci Iudeii I-au răspuns: “Nu spunem noi bine că ești samaritean și că ai un demon?”


\par }{\PP \VS{49}Isus a răspuns: “{\WJ{Eu nu am niciun demon, dar Eu Îl cinstesc pe Tatăl Meu, iar voi Mă necinstiți.
}}

\VS{50}{\WJ{Dar eu nu caut gloria mea. Există cineva care caută și judecă.
}}

\VS{51}{\WJ{Mai mult ca sigur, vă spun că, dacă o persoană respectă cuvântul meu, nu va vedea niciodată moartea.”
}}


\par }{\PP \VS{52}Atunci iudeii i-au zis: “Acum știm că ai un demon. Avraam a murit, ca și profeții, iar tu spui:
{\WJ{“Dacă un om păzește cuvântul Meu, nu va gusta niciodată moartea”.
}}

\VS{53}Ești tu mai mare decât tatăl nostru Avraam, care a murit? Profeții au murit. Cine te dai tu drept cine ești?”.


\par }{\PP \VS{54}Isus a răspuns: “{\WJ{Dacă mă slăvesc pe mine însumi, slava mea nu este nimic. Cel care mă slăvește pe mine este Tatăl meu, despre care voi spuneți că este Dumnezeul nostru.
}}

\VS{55}{\WJ{Voi nu L-ați cunoscut, dar Eu Îl cunosc. Dacă aș spune: “Nu-l cunosc”, aș fi ca voi, un mincinos. Dar eu îl cunosc și îi respect cuvântul.
}}

\VS{56}{\WJ{Tatăl vostru Avraam s-a bucurat să vadă ziua mea. A văzut-o și s-a bucurat”.
}}


\par }{\PP \VS{57}Și Iudeii i-au zis: “Nu ești încă în vârstă de cincizeci de ani! L-ai văzut tu pe Avraam?”


\par }{\PP \VS{58}Isus le-a zis: “{\WJ{Adevărat vă spun că, înainte ca Avraam să fi existat, EU SUNT.”
}}


\par }{\PP \VS{59}Și au luat pietre ca să arunce în El; dar Isus S-a ascuns și a ieșit din templu, trecând prin mijlocul lor, și a trecut pe lângă ei.



\par }\Chap{9}{\PP \VerseOne{1}Pe când trecea pe acolo, a văzut un om orb din naștere.

\VS{2}Ucenicii Lui l-au întrebat: “Rabi, cine a păcătuit, omul acesta sau părinții lui, de s-a născut orb?”


\par }{\PP \VS{3}Isus a răspuns: “{\WJ{Omul acesta n-a păcătuit, nici părinții lui, ci ca să se arate în el lucrările lui Dumnezeu.
}}

\VS{4}{\WJ{Eu trebuie să lucrez lucrările celui care m-a trimis, cât este ziuă. Se apropie noaptea, când nimeni nu mai poate lucra.
}}

\VS{5}{\WJ{Cât timp sunt în lume, eu sunt lumina lumii.”
}}

\VS{6}După ce a spus acestea, a scuipat pe pământ, a făcut noroi cu saliva, a uns cu noroi ochii orbului

\VS{7}și i-a spus:
{\WJ{“Du-te și te spală în piscina din Siloam” (care înseamnă “trimis”). El s-a dus, s-a spălat și s-a întors văzând.
}}


\par }{\PP \VS{8}De aceea, vecinii și cei ce văzuseră mai înainte că era orb, ziceau: “Nu este acesta cel ce ședea și cerșea?”

\VS{9}Alții spuneau: “El este”. Alții încă mai spuneau: “Seamănă cu el”.

\par }{\PP El a spus: “Eu sunt”.


\par }{\PP \VS{10}Și ei Îl întrebau: “Cum ți s-au deschis ochii?”


\par }{\PP \VS{11}El a răspuns: “Un om, numit Isus, a făcut noroi, mi-a uns ochii și mi-a zis:
{\WJ{“Du-te la iazul din Siloam și te spală. Așa că m-am dus, m-am spălat și am primit vederea.”
}}


\par }{\PP \VS{12}Atunci l-au întrebat: “Unde este?”

\par }{\PP El a spus: “Nu știu”.


\par }{\PP \VS{13}L-au dus pe cel orb la farisei.

\VS{14}Era Sabat când Isus a făcut noroiul și i-a deschis ochii.

\VS{15}De aceea și fariseii l-au întrebat iarăși cum și-a primit vederea. El le-a răspuns: “Mi-a pus noroi pe ochi, m-am spălat și văd”.


\par }{\PP \VS{16}Și unii dintre farisei ziceau: “Omul acesta nu este de la Dumnezeu, pentru că nu ține Sabatul.”

\par }{\PP Alții spuneau: “Cum poate un om păcătos să facă asemenea semne?” Astfel, a fost dezbinare între ei.


\par }{\PP \VS{17}De aceea au întrebat iarăși pe orb: “Ce zici despre El, pentru că ți-a deschis ochii?”

\par }{\PP El a spus: “Este un profet”.


\par }{\PP \VS{18}Iudeii n-au crezut despre El că fusese orb și că și-a recăpătat vederea, până ce au chemat pe părinții celui ce-și recăpătase vederea,

\VS{19}și i-au întrebat: “Acesta este fiul vostru, despre care spuneți că s-a născut orb? Cum deci vede acum?”


\par }{\PP \VS{20}Părinții lui le-au răspuns: “Știm că acesta este fiul nostru și că s-a născut orb;

\VS{21}dar nu știm cum vede acum, și nu știm cine i-a deschis ochii. El este în vârstă. Întrebați-l pe el. Va vorbi el însuși”.

\VS{22}Părinții lui au spus aceste lucruri pentru că se temeau de iudei; căci iudeii se înțeleseseră deja că, dacă cineva îl va recunoaște ca fiind Hristos, va fi dat afară din sinagogă.

\VS{23}De aceea, părinții lui au spus: “El este major. Întrebați-l”.


\par }{\PP \VS{24}Și au chemat a doua oară pe orbul acela și i-au zis: “Slăvește pe Dumnezeu. Știm că omul acesta este un păcătos.”


\par }{\PP \VS{25}El a răspuns: “Nu știu dacă este păcătos. Un lucru știu însă: că, deși eram orb, acum văd.”


\par }{\PP \VS{26}Și I-au zis iarăși: “Ce ți-a făcut? Cum ți-a deschis ochii?”


\par }{\PP \VS{27}El le-a răspuns: “V-am spus deja și nu m-ați ascultat. De ce vreți să auziți din nou? Nu vreți să deveniți și voi discipoli ai lui, nu-i așa?”


\par }{\PP \VS{28}Ei l-au insultat și au zis: “Tu ești ucenicul lui, iar noi suntem ucenicii lui Moise.

\VS{29}Noi știm că Dumnezeu a vorbit lui Moise. Dar în ceea ce-l privește pe acest om, nu știm de unde vine.”


\par }{\PP \VS{30}Omul le-a răspuns: “Ce minunat! Voi nu știți de unde vine, și totuși mi-a deschis ochii.

\VS{31}Știm că Dumnezeu nu ascultă de păcătoși, dar dacă cineva este închinător lui Dumnezeu și face voia Lui, El îl ascultă.

\VS{32}De la începutul lumii nu s-a auzit niciodată ca cineva să fi deschis ochii unui orb din naștere.

\VS{33}Dacă omul acesta nu era de la Dumnezeu, nu putea face nimic.”


\par }{\PP \VS{34}Ei I-au răspuns: “Tu, care te-ai născut cu totul în păcate, ne înveți pe noi?” Atunci l-au aruncat afară.


\par }{\PP \VS{35}Isus a auzit că l-au dat afară și, găsindu-l, a zis: “{\WJ{Crezi tu în Fiul lui Dumnezeu?”
}}


\par }{\PP \VS{36}El a răspuns: “Cine este El, Doamne, ca să cred în El?”


\par }{\PP \VS{37}Isus i-a zis: “{\WJ{L-ați văzut amândoi și El este Cel ce vorbește cu voi.”
}}


\par }{\PP \VS{38}El a zis: “Doamne, cred!” și s-a închinat lui.


\par }{\PP \VS{39}Isus a zis:
{\WJ{“Eu am venit în lumea aceasta ca să fac judecată, ca cei ce nu văd să vadă, iar cei ce văd să devină orbi.”
}}


\par }{\PP \VS{40}Fariseii care erau cu El, auzind aceste lucruri, I-au zis: “Suntem și noi orbi?”


\par }{\PP \VS{41}Isus le-a zis: “{\WJ{Dacă ați fi orbi, n-ați avea nici un păcat; dar acum ziceți: “Vedem”. Prin urmare, păcatul vostru rămâne.
}}



\par }\Chap{10}{\PP \VerseOne{1}“{\WJ{Adevărat vă spun că cel ce nu intră pe ușă în staulul oilor, ci se urcă pe altă cale, este un hoț și un tâlhar.
}}

\VS{2}{\WJ{Dar cel care intră pe ușă este păstorul oilor.
}}

\VS{3}{\WJ{Portarul îi deschide poarta, și oile ascultă de glasul lui. El își cheamă oile pe nume și le conduce afară.
}}

\VS{4}{\WJ{Ori de câte ori își scoate oile, el merge înaintea lor, și oile îl urmează, pentru că îi cunosc glasul.
}}

\VS{5}{\WJ{Ele nu vor urma nicidecum pe un străin, ci vor fugi de el, pentru că nu cunosc glasul străinilor.”
}}

\VS{6}Isus le-a spus această parabolă, dar ei nu au înțeles ce le spunea.


\par }{\PP \VS{7}Isus le-a zis din nou: “{\WJ{Adevărat, adevărat vă spun că Eu sunt ușa oilor.
}}

\VS{8}{\WJ{Toți cei care au venit înaintea mea sunt hoți și tâlhari, dar oile nu i-au ascultat.
}}

\VS{9}{\WJ{Eu sunt ușa. Dacă intră cineva prin mine, va fi mântuit, va intra și va ieși și va găsi pășune.
}}

\VS{10}{\WJ{Hoțul vine numai ca să fure, să ucidă și să distrugă. Eu am venit ca ei să aibă viață și să o aibă din belșug.
}}


\par }{\PP \VS{11}{\WJ{“Eu sunt păstorul cel bun. Păstorul cel bun își dă viața pentru oi.
}}

\VS{12}{\WJ{Cel care este angajat și nu păstor, care nu este proprietarul oilor, vede lupul venind, lasă oile și fuge. Lupul smulge oile și le împrăștie.
}}

\VS{13}{\WJ{Salariatului fuge, pentru că este salariat și nu are grijă de oi.
}}

\VS{14}{\WJ{Eu sunt păstorul cel bun. Eu îmi cunosc pe ale mele și sunt cunoscut de ale mele;
}}

\VS{15}{\WJ{așa cum Tatăl mă cunoaște pe mine și eu îl cunosc pe Tatăl. Eu îmi dau viața pentru oi.
}}

\VS{16}{\WJ{Am și alte oi care nu sunt din turma aceasta. Trebuie să le aduc și pe ele, și ele vor auzi glasul meu. Ele vor deveni o singură turmă cu un singur păstor.
}}

\VS{17}{\WJ{De aceea mă iubește Tatăl, pentru că îmi dau viața, ca să o iau din nou.
}}

\VS{18}{\WJ{Nimeni nu mi-o ia, ci eu mi-o pun singur. Eu am puterea să o dau și am puterea să o iau din nou. Am primit această poruncă de la Tatăl Meu.”
}}


\par }{\PP \VS{19}De aceea s-a făcut iarăși dezbinare între iudei din cauza acestor cuvinte.

\VS{20}Mulți dintre ei spuneau: “Are un demon și este nebun! De ce îl ascultați?”

\VS{21}Alții au spus: “Acestea nu sunt spusele unuia posedat de un demon. Nu este posibil ca un demon să deschidă ochii unui orb, nu-i așa?”


\par }{\PP \VS{22}Era Sărbătoarea Dedicării la Ierusalim.

\VS{23}Era iarnă, și Isus se plimba prin Templu, în pridvorul lui Solomon.

\VS{24}Iudeii s-au strâns în jurul lui și i-au zis: “Până când ne vei ține în suspans? Dacă tu ești Hristosul, spune-ne clar”.


\par }{\PP \VS{25}Isus le-a răspuns: “{\WJ{V-am spus, și nu credeți. Lucrările pe care le fac în numele Tatălui meu, acestea mărturisesc despre mine.
}}

\VS{26}{\WJ{Dar voi nu credeți, pentru că nu sunteți dintre oile Mele, așa cum v-am spus.}}

\VS{27}{\WJ{Oile Mele aud glasul Meu, și Eu le cunosc, și ele Mă urmează.
}}

\VS{28}{\WJ{Eu le dau viața veșnică. Ele nu vor pieri niciodată și nimeni nu le va smulge din mâna mea.
}}

\VS{29}{\WJ{Tatăl Meu, care Mi le-a dat mie, este mai mare decât toți. Nimeni nu poate să le smulgă din mâna Tatălui meu.
}}

\VS{30}{\WJ{Eu și Tatăl suntem una.”
}}


\par }{\PP \VS{31}De aceea iudeii au luat iarăși pietre ca să-L ucidă cu pietre.

\VS{32}Isus le-a răspuns: “{\WJ{Eu v-am arătat multe fapte bune de la Tatăl Meu. Pentru care dintre aceste fapte Mă ucideți cu pietre?”
}}


\par }{\PP \VS{33}Iudeii I-au răspuns: “Noi nu te ucidem cu pietre pentru o faptă bună, ci pentru blasfemie, pentru că tu, fiind om, te faci pe tine însuți Dumnezeu.”


\par }{\PP \VS{34}Isus le-a răspuns: “{\WJ{Nu este scris în legea voastră: “Eu am zis: “Voi sunteți dumnezei”?”}}

\VS{35}{\WJ{Dacă El i-a numit dumnezei pe cei la care a venit cuvântul lui Dumnezeu (și Scriptura nu poate fi călcată),
}}

\VS{36}{\WJ{spuneți voi despre Cel pe care Tatăl L-a sfințit și L-a trimis în lume: “Tu hulești”, pentru că Eu am zis: “Eu sunt Fiul lui Dumnezeu”?
}}

\VS{37}{\WJ{Dacă nu fac faptele Tatălui Meu, nu Mă credeți.
}}

\VS{38}{\WJ{Dar dacă le fac, chiar dacă nu mă credeți, credeți faptele, ca să știți și să credeți că Tatăl este în Mine și Eu în Tatăl.”
}}


\par }{\PP \VS{39}Ei au căutat din nou să-L prindă, dar El a scăpat din mâna lor.

\VS{40}S-a dus iarăși dincolo de Iordan, în locul unde botezase Ioan la început, și a rămas acolo.

\VS{41}Mulți veneau la el. Ei spuneau: “Într-adevăr, Ioan nu a făcut niciun semn, dar tot ce a spus Ioan despre omul acesta este adevărat.”

\VS{42}Mulți au crezut în el acolo.



\par }\Chap{11}{\PP \VerseOne{1}Un om bolnav, Lazăr, din Betania, din satul Mariei și al Martei, sora ei, era bolnav.

\VS{2}Maria, care unsese pe Domnul cu mir și Îi ștergea picioarele cu părul ei, era aceea al cărei frate, Lazăr, era bolnav.

\VS{3}De aceea surorile au trimis la el, spunând: “Doamne, iată că cel pentru care ai mare afecțiune este bolnav.”


\par }{\PP \VS{4}Dar Isus, auzind, a zis: “{\WJ{Boala aceasta nu este pentru moarte, ci pentru slava lui Dumnezeu, ca Fiul lui Dumnezeu să fie proslăvit prin ea.”
}}

\VS{5}Isus iubea pe Marta, pe sora ei și pe Lazăr.

\VS{6}De aceea, când a auzit că acesta era bolnav, a stat două zile în locul unde se afla.

\VS{7}După aceea, a zis ucenicilor:
{\WJ{“Să mergem din nou în Iudeea.”
}}


\par }{\PP \VS{8}Ucenicii L-au întrebat: “Rabi, iudeii tocmai voiau să Te ucidă cu pietre. Te duci din nou acolo?”


\par }{\PP \VS{9}Isus a răspuns: “{\WJ{Nu sunt oare douăsprezece ore de lumină? Dacă un om umblă ziua, nu se poticnește, pentru că vede lumina acestei lumi.
}}

\VS{10}{\WJ{Dar dacă un om umblă noaptea, se împiedică, pentru că lumina nu este în el.”
}}

\VS{11}El a spus aceste lucruri și, după aceea, le-a zis:
{\WJ{“Prietenul nostru Lazăr a adormit, dar eu mă duc ca să-l trezesc din somn.”
}}


\par }{\PP \VS{12}Și ucenicii au zis: “Doamne, dacă a adormit, își va reveni.”


\par }{\PP \VS{13}Isus vorbise despre moartea Sa, dar ei credeau că vorbea despre odihna în somn.

\VS{14}Atunci Isus le-a spus clar și răspicat: “{\WJ{Lazăr a murit.
}}

\VS{15}Mă
{\WJ{bucur, pentru voi, că nu am fost acolo, ca să credeți. Totuși, să mergem la el”.
}}


\par }{\PP \VS{16}Toma, zis Didim, a zis tovarășilor săi ucenici: “Să mergem și noi, ca să murim împreună cu El”.


\par }{\PP \VS{17}Când a venit Isus, a aflat că era deja de patru zile în mormânt.

\VS{18}Betania era aproape de Ierusalim, la o distanță de vreo cincisprezece stadii.

\VS{19}Mulți dintre iudei se alăturaseră femeilor din jurul Martei și al Mariei, ca să le consoleze cu privire la fratele lor.

\VS{20}Când a auzit Marta că vine Isus, s-a dus să-l întâmpine, dar Maria a rămas în casă.

\VS{21}Atunci Marta i-a zis lui Isus: “Doamne, dacă ai fi fost aici, fratele meu nu ar fi murit.

\VS{22}Chiar și acum știu că tot ce ceri de la Dumnezeu, Dumnezeu îți va da”.


\par }{\PP \VS{23}Isus i-a zis:
{\WJ{“Fratele tău va învia.”
}}


\par }{\PP \VS{24}Marta i-a zis: “Știu că va învia la înviere, în ziua de apoi.”


\par }{\PP \VS{25}Isus i-a zis: “{\WJ{Eu sunt învierea și viața. Cel ce crede în Mine va trăi, chiar dacă va muri.
}}

\VS{26}{\WJ{Oricine trăiește și crede în Mine nu va muri niciodată. Crezi tu asta?”
}}


\par }{\PP \VS{27}Ea i-a zis: “Da, Doamne. Am ajuns să cred că Tu ești Hristosul, Fiul lui Dumnezeu, Cel care vine în lume.”


\par }{\PP \VS{28}După ce a zis acestea, s-a dus și a chemat pe ascuns pe Maria, sora ei, zicând: “Învățătorul este aici și te cheamă.”


\par }{\PP \VS{29}Când a auzit aceasta, s-a sculat repede și s-a dus la el.

\VS{30}Isus nu intrase încă în sat, ci era în locul unde L-a întâlnit Marta.

\VS{31}Iudeii care erau cu ea în casă și o consolau, când au văzut-o pe Maria că s-a sculat repede și a ieșit, au urmărit-o și au zis: “Se duce la mormânt ca să plângă acolo.”


\par }{\PP \VS{32}Maria, când a ajuns unde era Isus și L-a văzut, a căzut la picioarele Lui și I-a zis: “Doamne, dacă ai fi fost aici, fratele meu n-ar fi murit.”


\par }{\PP \VS{33}Isus, văzând-o pe ea plângând și pe iudeii care plângeau împreună cu ea, a gemut în duh și s-a tulburat

\VS{34}și a zis:
{\WJ{“Unde L-ați pus?”
}}

\par }{\PP I-au spus: “Doamne, vino și vezi”.


\par }{\PP \VS{35}Isus a plâns.


\par }{\PP \VS{36}Iudeii ziceau deci: “Vedeți câtă dragoste avea pentru el!”

\VS{37}Unii dintre ei spuneau: “Nu putea oare acest om, care a deschis ochii orbului, să împiedice și pe acesta să moară?”


\par }{\PP \VS{38}Isus a venit la mormânt, gemând iarăși în sine. Era o peșteră, și o piatră era așezată împotriva ei.

\VS{39}Isus a zis:
{\WJ{“Scoateți piatra”.
}}

\par }{\PP Marta, sora celui mort, I-a zis: “Doamne, la ora aceasta este o putoare, căci este mort de patru zile.”


\par }{\PP \VS{40}Isus i-a zis: “{\WJ{Nu ți-am spus Eu că, dacă vei crede, vei vedea slava lui Dumnezeu?”
}}


\par }{\PP \VS{41}Și au îndepărtat piatra din locul unde zăcea mortul. Isus și-a ridicat ochii și a zis:
{\WJ{“Tată, îți mulțumesc că m-ai ascultat.
}}

\VS{42}{\WJ{Știu că întotdeauna mă asculți, dar, din cauza mulțimii care stătea în jur, am spus aceasta, ca să creadă că tu m-ai trimis.”
}}

\VS{43}După ce a spus aceasta, a strigat cu glas tare:
{\WJ{“Lazăr, ieși afară!”
}}


\par }{\PP \VS{44}Cel mort a ieșit, legat de mâini și de picioare, și avea fața înfășurată cu o pânză.

\par }{\PP Isus le-a zis:
{\WJ{“Eliberați-l și lăsați-l să plece”.
}}


\par }{\PP \VS{45}De aceea mulți dintre iudeii care au venit la Maria și au văzut ce făcea Isus au crezut în El.

\VS{46}Dar unii dintre ei s-au dus la farisei și le-au spus ce făcuse Isus.

\VS{47}Preoții cei mai de seamă și fariseii au adunat deci un consiliu și au zis: “Ce facem? Căci omul acesta face multe semne.

\VS{48}Dacă îl lăsăm așa, toată lumea va crede în el, iar romanii vor veni și ne vor lua atât locul nostru, cât și națiunea noastră.”


\par }{\PP \VS{49}Dar unul dintre ei, Caiafa, care era mare preot în anul acela, le-a zis: “Voi nu știți nimic,

\VS{50}și nici nu vă gândiți că ne este de folos ca un singur om să moară pentru popor și să nu piară tot neamul.”

\VS{51}Or, el nu a spus acest lucru de la sine, ci, fiind mare preot în acel an, a profețit că Isus va muri pentru națiune,

\VS{52}și nu numai pentru națiune, ci și pentru a-i aduna la un loc pe copiii lui Dumnezeu care sunt împrăștiați.

\VS{53}Astfel, din ziua aceea s-au sfătuit ca să îl omoare.

\VS{54}De aceea, Isus nu a mai umblat în mod deschis printre iudei, ci a plecat de acolo în ținutul de lângă pustiu, într-o cetate numită Efraim. Acolo a rămas cu discipolii Săi.


\par }{\PP \VS{55}Și se apropia Paștele iudeilor. Mulți se urcau din țară la Ierusalim înainte de Paște, ca să se purifice.

\VS{56}Apoi căutau pe Isus și vorbeau între ei, în timp ce stăteau în templu: “Ce credeți voi — că nu vine deloc la sărbătoare?”

\VS{57}Preoții cei mai de seamă și fariseii porunciseră ca, dacă cineva știa unde se află, să anunțe, ca să îl prindă.



\par }\Chap{12}{\PP \VerseOne{1}Cu șase zile înainte de Paști, Isus a venit la Betania, unde era Lazăr, care era mort și pe care l-a înviat din morți.

\VS{2}Și I-au pregătit acolo o cină. Marta a servit, dar Lazăr era unul dintre cei care stăteau la masă cu el.

\VS{3}Maria a luat un kilogram de mir de nard curat, foarte prețios, a uns picioarele lui Isus și I-a șters picioarele cu părul ei. Casa s-a umplut de mireasma unguentului.


\par }{\PP \VS{4}Atunci Iuda Iscarioteanul, fiul lui Simon, unul din ucenicii Lui, care avea să-L trădeze, a zis:

\VS{5}“De ce nu s-a vândut mirul acesta cu trei sute de denari și nu s-a dat săracilor?”

\VS{6}Și a spus aceasta nu pentru că-i păsa de săraci, ci pentru că era hoț și, având cutia cu bani, obișnuia să fure ce se punea în ea.


\par }{\PP \VS{7}Dar Isus a zis: “{\WJ{Las-o în pace. Ea a păstrat aceasta pentru ziua înmormântării mele.
}}

\VS{8}{\WJ{Căci voi aveți întotdeauna săraci cu voi, dar pe Mine nu Mă aveți întotdeauna.”
}}


\par }{\PP \VS{9}O mare mulțime de iudei au aflat că El era acolo și au venit, nu numai de dragul lui Isus, ci ca să vadă și pe Lazăr, pe care El îl înviase din morți.

\VS{10}Dar preoții cei mai de seamă au uneltit ca să îl omoare și pe Lazăr,

\VS{11}pentru că, din cauza lui, mulți dintre iudei au plecat și au crezut în Isus.


\par }{\PP \VS{12}A doua zi, o mare mulțime a venit la ospăț. Când au auzit că Isus vine la Ierusalim,

\VS{13}au luat ramuri de palmier, au ieșit în întâmpinarea lui și au strigat: “Osana! Binecuvântat este cel care vine în numele Domnului, Regele lui Israel!”.


\par }{\PP \VS{14}Isus a găsit un măgăruș tânăr și a șezut pe el. După cum este scris:

\VS{15}“Nu te teme, fiică a Sionului! Iată, Regele tău vine, șezând pe un mânz de măgăriță.”

\VS{16}La început, discipolii lui nu au înțeles aceste lucruri, dar când Isus a fost glorificat, atunci și-au adus aminte că aceste lucruri au fost scrise despre el și că ei i-au făcut aceste lucruri.

\VS{17}Așadar, mulțimea care era cu el când l-a chemat pe Lazăr din mormânt și l-a înviat din morți mărturisea despre aceasta.

\VS{18}Din această cauză și mulțimea a mers în întâmpinarea lui, pentru că auzise că el făcuse acest semn.

\VS{19}De aceea fariseii ziceau între ei: “Vedeți cum nu realizați nimic. Iată că lumea a mers după el”.


\par }{\PP \VS{20}Printre cei ce se suiseră să se închine la sărbătoare erau niște greci.

\VS{21}Aceștia au venit la Filip, care era din Betsaida Galileii, și l-au întrebat, zicând: “Domnule, vrem să-L vedem pe Isus.”

\VS{22}Filip a venit și i-a spus lui Andrei și, la rândul său, Andrei a venit cu Filip și i-au spus lui Isus.


\par }{\PP \VS{23}Isus le-a răspuns: “{\WJ{A venit vremea ca Fiul Omului să fie proslăvit.
}}

\VS{24}{\WJ{Adevărat vă spun că, dacă un bob de grâu nu cade în pământ și nu moare, el rămâne singur. Dar, dacă moare, face mult rod.
}}

\VS{25}{\WJ{Cine își iubește viața o va pierde. Cel care își urăște viața în lumea aceasta o va păstra pentru viața veșnică.
}}

\VS{26}{\WJ{Dacă cineva îmi slujește, să mă urmeze. Unde sunt eu, acolo va fi și slujitorul meu. Dacă cineva îmi slujește, Tatăl îl va cinsti.
}}


\par }{\PP \VS{27}{\WJ{“Acum sufletul meu este tulburat. Ce să spun? 'Tată, salvează-mă din acest timp?'. Dar am venit în acest timp pentru această cauză.
}}

\VS{28}{\WJ{Tată, slăvește numele Tău!”
}}

\par }{\PP Atunci un glas a ieșit din cer și a zis: “Am proslăvit-o și o voi proslăvi din nou”.


\par }{\PP \VS{29}De aceea mulțimea care stătea de față și auzea, zicea că a tunat. Alții spuneau: “Un înger i-a vorbit”.


\par }{\PP \VS{30}Isus a răspuns: “{\WJ{Glasul acesta n-a venit pentru Mine, ci pentru voi.
}}

\VS{31}{\WJ{Acum este judecata lumii acesteia. Acum prințul acestei lumi va fi alungat.
}}

\VS{32}{\WJ{Iar Eu, dacă voi fi ridicat de pe pământ, voi atrage la Mine toate popoarele.”
}}

\VS{33}Dar El a spus aceasta, însemnând prin ce fel de moarte va muri.


\par }{\PP \VS{34}Și mulțimea i-a răspuns: “Am auzit din Lege că Hristosul rămâne pentru totdeauna. Cum spuneți voi:
{\WJ{'Fiul Omului trebuie să fie înălțat'? Cine este acest Fiu al Omului?”
}}


\par }{\PP \VS{35}Isus le-a zis: “{\WJ{Încă puțină vreme lumina este cu voi. Umblați cât timp aveți lumina, ca să nu vă cuprindă întunericul. Cel care umblă în întuneric nu știe încotro se îndreaptă.
}}

\VS{36}{\WJ{Cât timp aveți lumina, credeți în lumină, ca să deveniți copii ai luminii.”}} Isus a spus aceste lucruri, apoi a plecat și s-a ascuns de ei.

\VS{37}Dar, deși făcuse atâtea semne înaintea lor, ei nu credeau în el,

\VS{38}pentru ca să se împlinească cuvântul proorocului Isaia, pe care îl rostise:

\par }{\Q “Doamne, cine a crezut raportul nostru?

\par }{\QB Cui s-a arătat brațul Domnului?”


\par }{\PP \VS{39}De aceea nu puteau să creadă, pentru că Isaia a spus din nou:


\par }{\Q \VS{40}“El le-a orbit ochii și le-a împietrit inima,

\par }{\QB ca nu cumva să vadă cu ochii lor,

\par }{\QB și să perceapă cu inima lor,

\par }{\QB și s-ar întoarce,

\par }{\QB și eu îi voi vindeca.”


\par }{\PP \VS{41}Isaia a spus aceste lucruri când a văzut slava Lui și a vorbit despre El.

\VS{42}Cu toate acestea, chiar și mulți dintre dregători au crezut în El, dar, din pricina fariseilor, n-au mărturisit-o, ca să nu fie scoși din sinagogă,

\VS{43}pentru că ei iubeau mai mult lauda oamenilor decât lauda lui Dumnezeu.


\par }{\PP \VS{44}Isus a strigat și a zis: “{\WJ{Cine crede în Mine, nu crede în Mine, ci în Cel ce M-a trimis pe Mine.
}}

\VS{45}{\WJ{Cine Mă vede pe Mine vede pe Cel ce M-a trimis pe Mine.
}}

\VS{46}{\WJ{Eu am venit ca o lumină în lume, pentru ca oricine crede în Mine să nu rămână în întuneric.
}}

\VS{47}{\WJ{Dacă cineva ascultă cele spuse de mine și nu crede, eu nu-l judec. Căci nu am venit să judec lumea, ci să salvez lumea.
}}

\VS{48}{\WJ{Cine mă respinge pe mine și nu primește spusele mele, are pe cineva care îl judecă. Cuvântul pe care l-am rostit îl va judeca în ziua de apoi.
}}

\VS{49}{\WJ{Căci nu am vorbit de la mine însumi, ci Tatăl care m-a trimis mi-a dat o poruncă, ce să spun și ce să vorbesc.
}}

\VS{50}{\WJ{Știu că porunca Lui este viața veșnică. Așadar, lucrurile pe care le spun, așa cum mi-a spus Tatăl, așa vorbesc și eu.”
}}



\par }\Chap{13}{\PP \VerseOne{1}Înainte de sărbătoarea Paștelui, Isus, știind că a venit vremea să plece din lumea aceasta la Tatăl, a iubit pe ai Săi care erau în lume și i-a iubit până la sfârșit.

\VS{2}În timpul cinei, când diavolul pusese deja în inima lui Iuda Iscarioteanul, fiul lui Simon, ca să-l trădeze,

\VS{3}Isus, știind că Tatăl îi dăduse toate lucrurile în mâinile sale și că el a venit de la Dumnezeu și că se duce la Dumnezeu,

\VS{4}s-a ridicat de la cină și și-a dat la o parte hainele de pe el. A luat un prosop și și-a înfășurat un prosop în jurul taliei.

\VS{5}Apoi a turnat apă în lighean și a început să spele picioarele discipolilor și să le șteargă cu prosopul care era înfășurat în jurul său.

\VS{6}Apoi a venit la Simon Petru. Acesta i-a zis: “Doamne, îmi speli tu mie picioarele?”.


\par }{\PP \VS{7}Isus i-a răspuns: “{\WJ{Nu știi ce fac acum, dar vei înțelege mai târziu.”
}}


\par }{\PP \VS{8}Petru i-a zis: “Nu-mi vei spăla niciodată picioarele!”

\par }{\PP Isus i-a răspuns:
{\WJ{“Dacă nu te spăl, nu ai parte cu Mine”.
}}


\par }{\PP \VS{9}Simon Petru I-a zis: “Doamne, nu numai picioarele mele, ci și mâinile și capul meu!”


\par }{\PP \VS{10}Isus i-a zis: “Cineva
{\WJ{care s-a scăldat nu are nevoie decât să-și spele picioarele, ci este cu totul curat. Voi sunteți curați, dar nu toți”.
}}

\VS{11}Căci îl cunoștea pe cel care avea să-l trădeze; de aceea a spus:
{\WJ{“Nu sunteți toți curați.”
}}

\VS{12}După ce le-a spălat picioarele, și-a pus la loc haina de pe el și a șezut din nou jos, le-a zis:
{\WJ{“Știți ce v-am făcut?
}}

\VS{13}{\WJ{Voi Mă numiți: “Învățătorule” și “Doamne”. Spuneți corect, căci așa sunt.
}}

\VS{14}Dacă
{\WJ{deci Eu, Domnul și Învățătorul, v-am spălat picioarele, și voi trebuie să vă spălați picioarele unii altora.
}}

\VS{15}{\WJ{Căci v-am dat o pildă, ca și voi să faceți cum v-am făcut eu.
}}

\VS{16}{\WJ{Adevărat vă spun că nu este un slujitor mai mare decât stăpânul său și nici cel trimis nu este mai mare decât cel care l-a trimis.
}}

\VS{17}{\WJ{Dacă știți aceste lucruri, ferice de voi dacă le veți face.
}}

\VS{18}{\WJ{Nu vorbesc cu privire la voi toți. Știu pe cine am ales, dar pentru ca să se împlinească Scriptura: “Cel care mănâncă pâine cu mine și-a ridicat călcâiul împotriva mea”.
}}

\VS{19}De
{\WJ{acum înainte vă spun înainte de a se întâmpla, pentru ca, atunci când se va întâmpla, să credeți că Eu sunt.
}}

\VS{20}Cu
{\WJ{siguranță vă spun: cine primește pe oricine îl trimit Eu, pe Mine Mă primește; și cine Mă primește pe Mine, Îl primește pe Cel care M-a trimis.”
}}


\par }{\PP \VS{21}După ce a zis acestea, Isus s-a tulburat în duhul său și a zis: “{\WJ{Adevărat vă spun că unul din voi Mă va trăda.”
}}


\par }{\PP \VS{22}Ucenicii se uitau unii la alții, și nu știau despre cine vorbea.

\VS{23}Unul dintre ucenicii Săi, pe care Isus îl iubea, era la masă și se sprijinea de pieptul lui Isus.

\VS{24}Simon Petru i-a făcut semn și i-a zis: “Spune-ne cine este acela despre care vorbește.”


\par }{\PP \VS{25}El, sprijinindu-se pe pieptul lui Isus, L-a întrebat: “Doamne, cine este?”


\par }{\PP \VS{26}Isus a răspuns: “{\WJ{Acela este cel căruia îi voi da această bucată de pâine, după ce o voi înmuia.” După ce a înmuiat bucata de pâine, a dat-o lui Iuda, fiul lui Simon Iscarioteanul.
}}

\VS{27}După bucățica de pâine, atunci Satana a intrat în el.

\par }{\PP Atunci Isus i-a zis: “{\WJ{Ce faci, fă repede!”.
}}


\par }{\PP \VS{28}Nimeni de la masă nu știa de ce îi spunea așa.

\VS{29}Fiindcă unii credeau, pentru că Iuda avea cutia cu bani, că Isus i-a spus: “Cumpără ce ne trebuie pentru sărbătoare”, sau că trebuie să dea ceva săracilor.

\VS{30}De aceea, după ce a primit bucățica aceea, a ieșit îndată. Era noapte.


\par }{\PP \VS{31}După ce a ieșit, Isus a zis: “{\WJ{Acum Fiul Omului a fost proslăvit și Dumnezeu a fost proslăvit în El.
}}

\VS{32}{\WJ{Dacă Dumnezeu a fost proslăvit în El, Dumnezeu Îl va proslăvi și în El însuși, și Îl va proslăvi îndată.
}}

\VS{33}{\WJ{Copilașilor, voi mai fi cu voi încă puțin timp. Mă veți căuta și, după cum am spus iudeilor: “Unde Mă duc Eu, voi nu puteți veni”, tot așa vă spun și acum.
}}

\VS{34}{\WJ{Vă dau o poruncă nouă: să vă iubiți unii pe alții. Așa cum v-am iubit Eu, așa să vă iubiți și voi unii pe alții.
}}

\VS{35}{\WJ{Prin aceasta va cunoaște toată lumea că sunteți discipolii mei, dacă vă veți iubi unii pe alții.”
}}


\par }{\PP \VS{36}Simon Petru I-a zis: “Doamne, unde Te duci?”

\par }{\PP Isus i-a răspuns: “{\WJ{Unde Mă duc Eu, nu poți să mă urmezi acum, dar mă vei urma după aceea.”
}}


\par }{\PP \VS{37}Petru I-a zis: “Doamne, de ce nu pot să Te urmez acum? Îmi voi da viața pentru Tine”.


\par }{\PP \VS{38}Isus i-a răspuns: “{\WJ{Vrei să-ți dai viața pentru Mine? Adevărat îți spun că nu va cânta cocoșul până când nu te vei lepăda de mine de trei ori.
}}



\par }\Chap{14}{\PP \VerseOne{1}{\WJ{“Să nu-ți tulburi inima. Credeți în Dumnezeu. Credeți și în mine.
}}

\VS{2}{\WJ{În casa Tatălui meu sunt multe case. Dacă nu ar fi fost așa, v-aș fi spus. Am să vă pregătesc un loc.
}}

\VS{3}{\WJ{Dacă Mă duc și vă pregătesc un loc, Mă voi întoarce și vă voi primi la Mine, pentru ca acolo unde sunt Eu, să fiți și voi.
}}

\VS{4}{\WJ{Voi știți unde mă duc și știți calea.”
}}


\par }{\PP \VS{5}Toma I-a zis: “Doamne, noi nu știm unde Te duci. Cum am putea cunoaște calea?”


\par }{\PP \VS{6}Isus i-a zis: “{\WJ{Eu sunt calea, adevărul și viața. Nimeni nu vine la Tatăl, decât prin mine.
}}

\VS{7}{\WJ{Dacă M-ai fi cunoscut pe Mine, L-ai fi cunoscut și pe Tatăl Meu. De acum înainte, tu îl cunoști și l-ai văzut.”
}}


\par }{\PP \VS{8}Filip i-a zis: “Doamne, arată-ne pe Tatăl și ne va fi de ajuns.”


\par }{\PP \VS{9}Isus i-a zis: “De
{\WJ{atâta vreme sunt cu tine și nu Mă cunoști, Filip? Cine M-a văzut pe Mine, L-a văzut pe Tatăl. Cum spui: “Arată-ne pe Tatăl?”
}}

\VS{10}{\WJ{Nu crezi că Eu sunt în Tatăl și Tatăl în mine? Cuvintele pe care vi le spun nu le spun de la mine însumi, ci Tatăl, care trăiește în mine, face lucrările Lui.
}}

\VS{11}{\WJ{Credeți-Mă că Eu sunt în Tatăl și Tatăl în Mine, sau credeți-Mă din cauza faptelor înseși.
}}

\VS{12}{\WJ{Adevărat vă spun că, cel ce crede în Mine, va face și el lucrările pe care le fac Eu; și va face lucrări mai mari decât acestea, pentru că Eu Mă duc la Tatăl Meu.
}}

\VS{13}{\WJ{Orice veți cere în numele Meu, voi face, pentru ca Tatăl să fie proslăvit în Fiul.
}}

\VS{14}{\WJ{Dacă veți cere ceva în numele Meu, Eu voi face.
}}

\VS{15}{\WJ{Dacă Mă iubiți, păziți poruncile Mele.
}}

\VS{16}{\WJ{Eu mă voi ruga Tatălui și el vă va da un alt Sfetnic, ca să fie cu voi în veci:
}}

\VS{17}{\WJ{Duhul adevărului, pe care lumea nu-l poate primi, pentru că nu-l vede și nu-l cunoaște. Voi îl cunoașteți, pentru că el trăiește cu voi și va fi în voi.
}}

\VS{18}{\WJ{Nu vă voi lăsa orfani. Eu voi veni la voi.
}}

\VS{19}{\WJ{Încă puțină vreme și lumea nu mă va mai vedea, dar voi mă veți vedea. Pentru că Eu trăiesc, veți trăi și voi.
}}

\VS{20}{\WJ{În ziua aceea, veți cunoaște că Eu sunt în Tatăl Meu, și voi în Mine, și Eu în voi.
}}

\VS{21}{\WJ{Cine are poruncile Mele și le păzește, acela este unul care Mă iubește. Cel care mă iubește pe mine va fi iubit de Tatăl meu, iar eu îl voi iubi și mă voi arăta lui.”
}}


\par }{\PP \VS{22}Iuda (nu Iscarioteanul) I-a zis: “Doamne, ce s-a întâmplat de Te-ai arătat nouă și nu lumii?”


\par }{\PP \VS{23}Isus i-a răspuns: “{\WJ{Dacă Mă iubește cineva, va păzi cuvântul Meu. Tatăl Meu îl va iubi și noi vom veni la el și ne vom face casa la el.
}}

\VS{24}{\WJ{Cine nu mă iubește pe mine nu păzește cuvintele mele. Cuvântul pe care îl auziți nu este al meu, ci al Tatălui care m-a trimis.
}}


\par }{\PP \VS{25}“{\WJ{V-am spus aceste lucruri în timp ce trăiam cu voi.
}}

\VS{26}{\WJ{Dar Consilierul, Duhul Sfânt, pe care Tatăl îl va trimite în numele Meu, vă va învăța toate lucrurile și vă va aminti tot ce v-am spus.
}}

\VS{27}{\WJ{Vă las pacea. Pacea Mea v-o dau vouă; nu cum dă lumea, vă dau Eu vouă. Să nu vă tulburați inima și să nu vă fie teamă.
}}

\VS{28}Ați
{\WJ{auzit cum v-am spus: “Mă duc și mă voi întoarce la voi”. Dacă m-ați fi iubit, v-ați fi bucurat, pentru că am spus: 'Mă duc la Tatăl meu'; căci Tatăl este mai mare decât mine.
}}

\VS{29}{\WJ{Acum v-am spus înainte de a se întâmpla, pentru ca, atunci când se va întâmpla, să credeți.
}}

\VS{30}{\WJ{Nu voi mai vorbi mult cu voi, pentru că vine prințul lumii și el nu are nimic în mine.
}}

\VS{31}{\WJ{Ci pentru ca lumea să știe că Eu iubesc pe Tatăl și că așa cum Mi-a poruncit Tatăl, așa fac și Eu. Ridicați-vă, să mergem de aici.
}}



\par }\Chap{15}{\PP \VerseOne{1}{\WJ{“Eu sunt vița cea adevărată și Tatăl Meu este viticultorul.
}}

\VS{2}{\WJ{Orice ramură din mine care nu face rod, El o taie. Orice ramură care dă rod, el o curăță, ca să dea mai mult rod.
}}

\VS{3}{\WJ{Voi sunteți deja tăiați curat, din cauza cuvântului pe care vi l-am spus.
}}

\VS{4}{\WJ{Rămâneți în mine și eu în voi. Așa cum ramura nu poate da rod de una singură dacă nu rămâne în viță, tot așa nici voi nu puteți, dacă nu rămâneți în mine.
}}

\VS{5}{\WJ{Eu sunt vița de vie. Voi sunteți ramurile. Cel care rămâne în mine și eu în el face multă roadă, pentru că în afară de mine nu puteți face nimic.
}}

\VS{6}{\WJ{Dacă un om nu rămâne în mine, este aruncat ca o ramură și se usucă; le adună, le aruncă în foc și sunt arse.
}}

\VS{7}{\WJ{Dacă rămâneți în mine și cuvintele mele rămân în voi, veți cere tot ce doriți și vi se va face.
}}


\par }{\PP \VS{8}{\WJ{Întru aceasta este slăvit Tatăl Meu, că voi faceți multă roadă, și astfel veți fi ucenicii Mei.
}}

\VS{9}{\WJ{Așa cum M-a iubit pe Mine Tatăl, așa v-am iubit și Eu pe voi. Rămâneți în dragostea mea.
}}

\VS{10}{\WJ{Dacă veți păzi poruncile Mele, veți rămâne în dragostea Mea, după cum și Eu am păzit poruncile Tatălui Meu și rămân în dragostea Lui.
}}

\VS{11}V-am spus
{\WJ{aceste lucruri pentru ca bucuria mea să rămână în voi și bucuria voastră să fie deplină.
}}


\par }{\PP \VS{12}“{\WJ{Aceasta este porunca Mea: Să vă iubiți unii pe alții, cum v-am iubit Eu pe voi.
}}

\VS{13}{\WJ{Nimeni nu are o dragoste mai mare decât aceasta: să-și dea cineva viața pentru prietenii săi.
}}

\VS{14}{\WJ{Voi sunteți prietenii Mei, dacă faceți tot ce vă poruncesc eu.
}}

\VS{15}{\WJ{Nu vă mai numesc robi, căci robul nu știe ce face stăpânul său. Ci v-am numit prieteni, pentru că tot ce am auzit de la Tatăl meu v-am făcut cunoscut.
}}

\VS{16}{\WJ{Nu voi m-ați ales pe mine, ci eu v-am ales pe voi și v-am rânduit ca să mergeți și să aduceți roade, iar roadele voastre să rămână, pentru ca tot ce veți cere de la Tatăl în numele meu, să vi se dea.
}}


\par }{\PP \VS{17}{\WJ{Vă poruncesc aceste lucruri, ca să vă iubiți unii pe alții.
}}

\VS{18}{\WJ{Dacă lumea vă urăște, să știți că și pe Mine M-a urât înainte de a vă urî pe voi.
}}

\VS{19}{\WJ{Dacă ați fi din lume, lumea și-ar iubi pe ai ei. Dar pentru că nu sunteți din lume, întrucât v-am ales din lume, de aceea vă urăște lumea.
}}

\VS{20}Aduceți-vă
{\WJ{aminte de cuvântul pe care vi l-am spus: 'Un rob nu este mai mare decât stăpânul său'. Dacă m-au persecutat pe mine, vă vor persecuta și pe voi. Dacă ei au păzit cuvântul meu, îl vor păzi și pe al vostru.
}}

\VS{21}{\WJ{Dar vă vor face toate acestea din cauza numelui meu, pentru că nu-l cunosc pe cel care m-a trimis.
}}

\VS{22}{\WJ{Dacă nu aș fi venit și nu le-aș fi vorbit, nu ar fi avut păcat; dar acum nu au nicio scuză pentru păcatul lor.
}}

\VS{23}{\WJ{Cine mă urăște pe mine, îl urăște și pe Tatăl meu.
}}

\VS{24}{\WJ{Dacă n-aș fi făcut între ei faptele pe care nimeni altcineva nu le-a făcut, ei n-ar fi avut păcat. Dar acum au văzut și m-au urât și pe mine și pe Tatăl meu.
}}

\VS{25}{\WJ{Dar aceasta s-a întâmplat pentru ca să se împlinească cuvântul care era scris în legea lor: “M-au urât fără motiv”.
}}


\par }{\PP \VS{26}{\WJ{“Când va veni Consilierul, pe care vi-L voi trimite de la Tatăl, Duhul adevărului, care purcede de la Tatăl, El va mărturisi despre Mine.
}}

\VS{27}{\WJ{Și voiveți mărturisi, pentru că ați fost cu mine de la început.
}}



\par }\Chap{16}{\PP \VerseOne{1}“{\WJ{V-am spus aceste lucruri ca să nu vă poticniți.
}}

\VS{2}{\WJ{Ei vă vor da afară din sinagogi. Da, vine vremea când oricine vă va ucide va crede că îi oferă un serviciu lui Dumnezeu.
}}

\VS{3}{\WJ{Ei vor face aceste lucruri pentru că nu l-au cunoscut pe Tatăl și nici pe mine.
}}

\VS{4}{\WJ{Dar v-am spus aceste lucruri pentru ca, atunci când va veni vremea, să vă amintiți că v-am vorbit despre ele. Nu v-am spus aceste lucruri de la început, pentru că eram cu voi.
}}

\VS{5}{\WJ{Dar acum mă duc la cel care m-a trimis și niciunul dintre voi nu mă întreabă: “Unde te duci?”
}}

\VS{6}{\WJ{Ci, pentru că v-am spus aceste lucruri, durerea v-a umplut inima.
}}

\VS{7}Cu toate
{\WJ{acestea, vă spun adevărul: este spre folosul vostru să plec, căci dacă nu plec, Consilierul nu va veni la voi. Dar dacă plec, vi-l voi trimite.
}}

\VS{8}{\WJ{Când va veni, el va convinge lumea despre păcat, despre dreptate și despre judecată;
}}

\VS{9}{\WJ{despre păcat, pentru că nu cred în mine;}}

\VS{10}despre
{\WJ{dreptate, pentru că mă duc la Tatăl meu și nu mă veți mai vedea;
}}

\VS{11}{\WJ{despre judecată, pentru că prințul acestei lumi a fost judecat.
}}


\par }{\PP \VS{12}{\WJ{Mai am încă multe lucruri să-ți spun, dar acum nu le poți suferi.
}}

\VS{13}{\WJ{Dar când va veni El, Duhul adevărului, vă va călăuzi în tot adevărul, căci nu va vorbi de la sine, ci va spune tot ce va auzi. El vă va vesti lucrurile care vor veni.
}}

\VS{14}{\WJ{El mă va proslăvi, pentru că va lua din ceea ce este al meu și vă va vesti.
}}

\VS{15}{\WJ{Toate lucrurile pe care le are Tatăl sunt ale mele; de aceea am spus că el ia din ale mele și vi le va vesti vouă.
}}


\par }{\PP \VS{16}“Încă
{\WJ{puțină vreme și nu Mă veți mai vedea. Din nou puțină vreme, și Mă veți vedea.”
}}


\par }{\PP \VS{17}De aceea unii dintre ucenicii Lui ziceau unii către alții: “Ce este aceasta, că ne spune: “{\WJ{Peste puțină vreme nu Mă veți vedea; și iarăși peste puțină vreme Mă veți vedea}}” și:
{\WJ{“Pentru că Mă duc la Tatăl”}}?”

\VS{18}Ei au spus deci: “Ce este aceasta pe care o spune:
{\WJ{“Puțin timp”? Nu știm ce vrea să spună”.
}}


\par }{\PP \VS{19}Isus a înțeles că voiau să-L întrebe și le-a zis: “Vă întrebați
{\WJ{între voi despre aceasta, că am zis: “Peste puțină vreme nu Mă veți vedea și iarăși peste puțină vreme Mă veți vedea?”
}}

\VS{20}{\WJ{Adevărat vă spun că voi veți plânge și vă veți tângui, iar lumea se va bucura. Voi veți fi întristați, dar întristarea voastră se va transforma în bucurie.
}}

\VS{21}{\WJ{Femeia, când naște, are tristețe, pentru că i-a sosit vremea. Dar, după ce a născut, nu-și mai amintește de chin, pentru bucuria că s-a născut o ființă umană în lume.
}}

\VS{22}{\WJ{De aceea, acum aveți întristare, dar vă voi revedea, și inima voastră se va bucura, și nimeni nu vă va lua bucuria.
}}


\par }{\PP \VS{23}{\WJ{În ziua aceea nu-mi veți pune întrebări. Adevărat vă spun că orice veți cere de la Tatăl în numele Meu, El vă va da.
}}

\VS{24}{\WJ{Până acum, nu ați cerut nimic în numele Meu. Cereți și veți primi, pentru ca bucuria voastră să fie deplină.
}}


\par }{\PP \VS{25}“{\WJ{V-am spus aceste lucruri prin figuri de stil. Dar vine vremea când nu vă voi mai vorbi cu figuri de stil, ci vă voi vorbi în mod clar despre Tatăl.
}}

\VS{26}{\WJ{În ziua aceea, veți cere în numele meu; și nu vă spun că eu mă voi ruga Tatălui pentru voi,
}}

\VS{27}{\WJ{pentru că însuși Tatăl vă iubește, pentru că voi m-ați iubit pe mine și ați crezut că am venit de la Dumnezeu.
}}

\VS{28}{\WJ{Eu am venit de la Tatăl și am venit în lume. Din nou, las lumea și mă duc la Tatăl.”
}}


\par }{\PP \VS{29}Ucenicii Lui I-au zis: “Iată că acum vorbești limpede și fără figuri de stil.

\VS{30}Acum știm că Tu știi toate lucrurile și nu este nevoie ca cineva să Te întrebe. Prin aceasta credem că ai venit de la Dumnezeu.”


\par }{\PP \VS{31}Isus le-a răspuns: “{\WJ{Credeți voi acum?
}}

\VS{32}{\WJ{Iată că vine vremea, și a venit, că veți fi împrăștiați, fiecare la locul lui, și Mă veți lăsa singur. Totuși, Eu nu sunt singur, pentru că Tatăl este cu Mine.
}}

\VS{33}{\WJ{V-am spus aceste lucruri pentru ca în Mine să aveți pace. În lume aveți necazuri; dar înveseliți-vă! Eu am biruit lumea”.
}}



\par }\Chap{17}{\PP \VerseOne{1}Isus a spus aceste lucruri, apoi, ridicându-Și ochii spre cer, a zis: “{\WJ{Tată, a venit vremea. Slăvește-l pe Fiul tău, ca și Fiul tău să te slăvească pe tine;
}}

\VS{2}{\WJ{așa cum i-ai dat autoritate peste toată făptura, așa va da viață veșnică tuturor celor pe care i-ai dat-o.
}}

\VS{3}{\WJ{Aceasta este viața veșnică: să te cunoască pe tine, singurul Dumnezeu adevărat, și pe cel pe care l-ai trimis, Isus Hristos.
}}

\VS{4}{\WJ{Eu te-am slăvit pe pământ. Am împlinit lucrarea pe care mi-ai dat-o să o fac.
}}

\VS{5}{\WJ{Acum, Tată, proslăvește-mă pe mine însumi cu gloria pe care o aveam cu tine înainte de a exista lumea.
}}


\par }{\PP \VS{6}“{\WJ{Eu am arătat Numele Tău poporului pe care mi l-ai dat din lume. Ei erau ai tăi și mi i-ai dat. Ei au respectat cuvântul tău.
}}

\VS{7}{\WJ{Acum au cunoscut că tot ce mi-ai dat vine de la Tine,
}}

\VS{8}{\WJ{pentru că le-am dat cuvintele pe care mi le-ai dat; și ei le-au primit și au știut cu siguranță că vin de la Tine. Ei au crezut că tu m-ai trimis.
}}

\VS{9}{\WJ{Eu mă rog pentru ei. Nu mă rog pentru lume, ci pentru cei pe care mi i-ai dat, pentru că ei sunt ai tăi.
}}

\VS{10}{\WJ{Toate lucrurile care sunt ale mele sunt ale tale, iar ale tale sunt ale mele și eu sunt glorificat în ele.
}}

\VS{11}{\WJ{Eu nu mai sunt în lume, dar aceștia sunt în lume, iar eu vin la voi. Tată sfânt, păzește-i prin numele tău, pe care mi l-ai dat mie, pentru ca ei să fie una, așa cum suntem și noi.
}}

\VS{12}{\WJ{Cât timp am fost cu ei în lume, i-am păzit în numele Tău. Am păzit pe cei pe care mi i-ai dat. Nici unul dintre ei nu este pierdut, cu excepția fiului pierzării, pentru ca Scriptura să se împlinească.
}}

\VS{13}{\WJ{Dar acum am venit la voi și spun aceste lucruri în lume, pentru ca ei să aibă bucuria mea deplină în ei înșiși.
}}

\VS{14}{\WJ{Le-am dat cuvântul Tău. Lumea i-a urât, pentru că nu sunt din lume, după cum nici Eu nu sunt din lume.
}}

\VS{15}{\WJ{Nu te rog să îi iei din lume, ci să îi păzești de cel rău.
}}

\VS{16}{\WJ{Ei nu sunt din lume, după cum nici Eu nu sunt din lume.
}}

\VS{17}{\WJ{Sfințește-i în adevărul Tău. Cuvântul Tău este adevărul.
}}

\VS{18}{\WJ{După cum m-ai trimis pe mine în lume, tot așa i-am trimis și eu pe ei în lume.
}}

\VS{19}De
{\WJ{dragul lor mă sfințesc pe mine însumi, pentru ca și ei înșiși să fie sfințiți în adevăr.
}}


\par }{\PP \VS{20}{\WJ{“Nu numai pentru aceștia mă rog, ci și pentru cei ce vor crede în Mine prin cuvântul lor,
}}

\VS{21}ca
{\WJ{toți să fie una, după cum Tu, Tată, ești în Mine și Eu în Tine, ca și ei să fie una în noi, pentru ca lumea să creadă că Tu M-ai trimis.
}}

\VS{22}{\WJ{Slava pe care Mi-ai dat-o Tu, Eu le-am dat-o lor, pentru ca ei să fie una, așa cum și noi suntem una,
}}

\VS{23}{\WJ{Eu în ei și Tu în Mine, ca ei să fie desăvârșiți într-o singură ființă, pentru ca lumea să știe că Tu M-ai trimis și că i-ai iubit, așa cum M-ai iubit pe Mine.
}}

\VS{24}{\WJ{Tată, doresc ca și aceia pe care mi i-ai dat să fie cu mine acolo unde sunt eu, ca să vadă slava pe care mi-ai dat-o, pentru că m-ai iubit înainte de întemeierea lumii.
}}

\VS{25}{\WJ{Tată dreptcredincios, lumea nu Te-a cunoscut, dar Eu Te-am cunoscut pe Tine; și aceștia au cunoscut că Tu M-ai trimis.
}}

\VS{26}{\WJ{Eule-am făcut cunoscut numele Tău și îl voi face cunoscut, pentru ca dragostea cu care M-ai iubit să fie în ei și Eu în ei.”
}}



\par }\Chap{18}{\PP \VerseOne{1}După ce a rostit aceste cuvinte, Isus a ieșit cu ucenicii Săi peste pârâul Chedron, unde era o grădină, în care a intrat împreună cu ucenicii Săi.

\VS{2}Iuda, cel care l-a trădat, cunoștea și el locul, pentru că Isus se întâlnea adesea acolo cu ucenicii săi.

\VS{3}Atunci Iuda, după ce a luat un detașament de soldați și de ofițeri de la preoții cei mai de seamă și de la farisei, a venit acolo cu felinare, torțe și arme.

\VS{4}Isus, deci, știind toate lucrurile care i se întâmplau, a ieșit și le-a zis:
{\WJ{“Pe cine căutați?”
}}


\par }{\PP \VS{5}Ei I-au răspuns: “Isus din Nazaret.”

\par }{\PP Isus le-a spus:
{\WJ{“Eu sunt.”
}}

\par }{\PP Și Iuda, care l-a trădat, stătea cu ei.

\VS{6}De aceea, când le-a zis:
{\WJ{“Eu sunt Acela”, ei s-au dat înapoi și au căzut la pământ.
}}


\par }{\PP \VS{7}De aceea i-a întrebat din nou: “Pe
{\WJ{cine căutați?”
}}

\par }{\PP Ei au spus: “Isus din Nazaret”.


\par }{\PP \VS{8}Isus a răspuns: “{\WJ{V-am spus că Eu sunt. De aceea, dacă Mă căutați pe Mine, lăsați-i pe aceștia să meargă pe drumul lor”
}}

\VS{9}ca să se împlinească cuvântul pe care l-a rostit: “{\WJ{Dintre cei pe care Mi i-ai dat, nu am pierdut pe niciunul”.
}}


\par }{\PP \VS{10}Simon Petru, având o sabie, a scos-o, a lovit pe robul marelui preot și i-a tăiat urechea dreaptă. Numele slujitorului era Malchus.

\VS{11}Isus i-a zis lui Petru:
{\WJ{“Pune sabia în teaca ei. Paharul pe care Mi l-a dat Tatăl, nu-l voi bea oare cu siguranță?”.
}}


\par }{\PP \VS{12}Atunci detașamentul, comandantul și ofițerii iudeilor au prins pe Isus și L-au legat,

\VS{13}și L-au dus mai întâi la Ana, care era socrul lui Caiafa, care era mare preot în anul acela.

\VS{14}Caiafa fusese cel care i-a sfătuit pe iudei că este bine ca un om să piară pentru popor.


\par }{\PP \VS{15}Simon Petru și un alt ucenic au urmat pe Isus. Acel ucenic era cunoscut de marele preot și a intrat cu Isus în curtea marelui preot;

\VS{16}dar Petru stătea afară la ușă. Deci celălalt ucenic, care era cunoscut de marele preot, a ieșit și a vorbit cu cea care păzea ușa și l-a adus pe Petru înăuntru.

\VS{17}Atunci servitoarea care păzea ușa i-a zis lui Petru: “Și tu ești unul dintre discipolii acestui om?”

\par }{\PP El a spus: “Nu sunt.”


\par }{\PP \VS{18}Și stăteau acolo slujitorii și ofițerii și făcuseră un foc de cărbuni, pentru că era frig. Ei se încălzeau. Petru era cu ei, stând în picioare și încălzindu-se.


\par }{\PP \VS{19}Marele preot a întrebat pe Isus despre ucenicii Săi și despre învățătura Lui.


\par }{\PP \VS{20}Isus i-a răspuns: “{\WJ{Eu am vorbit deschis lumii. Întotdeauna am învățat în sinagogi și în templu, unde se întâlnesc mereu evreii. Nu am spus nimic în secret.
}}

\VS{21}{\WJ{De ce mă întrebi pe mine? Întreabă-i pe cei care m-au ascultat ce le-am spus. Iată, ei știu ce am spus.”
}}


\par }{\PP \VS{22}După ce a zis acestea, unul din ofițerii care stăteau acolo a lovit pe Isus cu mâna, zicând: “Așa răspunzi Tu marelui preot?”


\par }{\PP \VS{23}Isus i-a răspuns: “Dacă
{\WJ{am vorbit de rău, mărturisește răul; dar dacă am spus bine, de ce mă bați?”
}}


\par }{\PP \VS{24}Și Ana l-a trimis legat la Caiafa, marele preot.


\par }{\PP \VS{25}Simon Petru stătea în picioare și se încălzea. Și ei i-au zis: “Nu cumva ești și tu unul dintre ucenicii Lui?”

\par }{\PP El a negat și a spus: “Nu sunt”.


\par }{\PP \VS{26}Unul din slujitorii marelui preot, rudă cu cel căruia Petru îi tăiase urechea, a zis: “Nu te-am văzut în grădină cu el?”


\par }{\PP \VS{27}Petru a negat din nou, și îndată a cântat cocoșul.


\par }{\PP \VS{28}Și au dus pe Isus de la Caiafa în pretoriu. Era devreme, și ei înșiși nu au intrat în pretoriu, ca să nu se spurce, ci să mănânce Paștele.

\VS{29}Pilat a ieșit deci la ei și i-a întrebat: “Ce acuzație aduceți împotriva omului acesta?”


\par }{\PP \VS{30}Ei i-au răspuns: “Dacă n-ar fi fost omul acesta un răufăcător, nu ți l-am fi dat ție.”


\par }{\PP \VS{31}Atunci Pilat le-a zis: “Luați-l voi înșivă și judecați-l după legea voastră.”

\par }{\PP De aceea Iudeii I-au zis: “Nu ne este îngăduit să omorâm pe nimeni”,

\VS{32}pentru ca să se împlinească cuvântul lui Isus, pe care-l spusese și care însemna prin ce fel de moarte avea să moară.


\par }{\PP \VS{33}Pilat a intrat din nou în pretoriu, a chemat pe Isus și I-a zis: “Tu ești regele iudeilor?”


\par }{\PP \VS{34}Isus i-a răspuns: “{\WJ{Spui lucrul acesta de unul singur, sau ți-au spus alții despre Mine?”
}}


\par }{\PP \VS{35}Pilat a răspuns: “Eu nu sunt iudeu, nu-i așa? Neamul tău și preoții cei mai de seamă te-au predat mie. Ce ai făcut?”


\par }{\PP \VS{36}Isus a răspuns: “{\WJ{Împărăția Mea nu este din lumea aceasta. Dacă Împărăția mea ar fi din lumea aceasta, atunci slujitorii mei s-ar lupta, ca să nu fiu predat iudeilor. Dar acum Împărăția mea nu este de aici.”
}}


\par }{\PP \VS{37}Și Pilat I-a zis: “Deci ești tu rege?”

\par }{\PP Isus a răspuns: “{\WJ{Tu spui că sunt rege. De aceea M-am născut și de aceea am venit în lume: ca să mărturisesc adevărul. Oricine este din adevăr ascultă glasul meu”.
}}


\par }{\PP \VS{38}Pilat I-a zis: “Ce este adevărul?”

\par }{\PP După ce a spus aceasta, a ieșit din nou la iudei și le-a zis: “Nu găsesc nici un temei pentru o acuzație împotriva lui.

\VS{39}Dar voi aveți obiceiul ca eu să vă eliberez pe cineva la Paște. Prin urmare, vreți să vi-l eliberez pe regele iudeilor?”


\par }{\PP \VS{40}Atunci toți au strigat iarăși: “Nu pe acesta, ci pe Baraba!” Or, Baraba era un tâlhar.



\par }\Chap{19}{\PP \VerseOne{1}Atunci Pilat a luat pe Isus și L-a biciuit.

\VS{2}Soldații au răsucit spini pentru a face o coroană, i-au pus-o pe cap și l-au îmbrăcat într-o haină de purpură.

\VS{3}Ei tot ziceau: “Slavă ție, rege al iudeilor!” și îl tot pălmuiau.


\par }{\PP \VS{4}Atunci Pilat a ieșit iarăși și le-a zis: “Iată, vi-l scot afară, ca să știți că nu găsesc nici un temei de acuzare împotriva lui.”


\par }{\PP \VS{5}Isus a ieșit afară, purtând cununa de spini și haina de purpură. Pilat le-a zis: “Iată-l pe Omul acesta!”.


\par }{\PP \VS{6}Atunci preoții cei mai de seamă și arhiereii, văzându-L, au strigat și au zis: “Răstignește-L! Răstignește!”

\par }{\PP Pilat le-a zis: “Luați-L voi înșivă și răstigniți-L, căci nu găsesc nici un temei de acuzare împotriva Lui.”


\par }{\PP \VS{7}Iudeii I-au răspuns: “Noi avem o lege și, după legea noastră, trebuie să moară, pentru că S-a făcut pe Sine Fiul lui Dumnezeu.”


\par }{\PP \VS{8}Când a auzit Pilat aceste cuvinte, s-a înspăimântat și mai mult.

\VS{9}A intrat din nou în pretoriu și a întrebat pe Isus: “De unde ești?” Dar Isus nu i-a dat niciun răspuns.

\VS{10}Atunci Pilat i-a zis: “Nu vorbești cu mine? Nu știi că am puterea să te eliberez și am puterea să te răstignesc?”


\par }{\PP \VS{11}Isus a răspuns: “{\WJ{N-aveți nici o putere împotriva Mea, dacă nu v-ar fi dată de sus. De aceea, cel care m-a predat vouă are un păcat mai mare.”
}}


\par }{\PP \VS{12}Atunci Pilat a vrut să-l elibereze, dar iudeii au strigat: “Dacă-l eliberezi pe omul acesta, nu ești prietenul Cezarului! Oricine se face rege vorbește împotriva Cezarului!”


\par }{\PP \VS{13}Pilat, auzind aceste cuvinte, a scos pe Isus afară și a șezut pe scaunul de judecată, într-un loc numit “Pavecernița”, dar în ebraică “Gabbatha”.

\VS{14}Era ziua pregătirii Paștelui, cam pe la ora șase. El le-a spus iudeilor: “Iată Regele vostru!”.


\par }{\PP \VS{15}Și au strigat: “Departe de El! Departe de el! Răstigniți-l!”

\par }{\PP Pilat i-a întrebat: “Să-l răstignesc pe regele vostru?”

\par }{\PP Preoții cei mai de seamă au răspuns: “Noi nu avem alt rege decât pe Cezar!”


\par }{\PP \VS{16}Și L-a dat să fie răstignit. Și au luat pe Isus și L-au dus.

\VS{17}El a ieșit, purtându-și crucea, la locul numit “Locul craniului”, care în ebraică se numește “Golgota”,

\VS{18}unde L-au răstignit și, împreună cu El, pe alți doi, de o parte și de alta, unul de fiecare parte, și pe Isus în mijloc.

\VS{19}Pilat a scris și un titlu și l-a pus pe cruce. Acolo era scris: “IISUS DIN NAZARETH, REGELE EVREILOR”.

\VS{20}De aceea, mulți dintre iudei citeau acest titlu, căci locul unde a fost răstignit Isus era aproape de cetate; și era scris în ebraică, în latină și în greacă.

\VS{21}De aceea, preoții cei mai de seamă ai iudeilor i-au spus lui Pilat: “Nu scrie: “Regele iudeilor”, ci: “El a spus: “Eu sunt Regele iudeilor”.””


\par }{\PP \VS{22}Pilat a răspuns: “Ceea ce am scris, am scris.”


\par }{\PP \VS{23}După ce au răstignit pe Isus, ostașii au luat hainele Lui și au făcut patru părți, câte o parte pentru fiecare ostaș, și tunica. Tunica era fără cusătură, țesută de sus în jos, de la un capăt la altul.

\VS{24}Și au zis unul către altul: “Să nu o rupem, ci să tragem la sorți pentru ea, ca să hotărâm a cui va fi”, ca să se împlinească Scriptura care zice

\par }{\Q “Au împărțit hainele mele între ei.

\par }{\QB Au tras la sorți pentru hainele mele.”

\par }{\MM De aceea soldații au făcut aceste lucruri.


\par }{\PP \VS{25}Iar lângă crucea lui Isus stăteau mama Lui, sora mamei Lui, Maria, nevasta lui Clopaș, și Maria Magdalena.

\VS{26}Când a văzut Isus pe mama Sa și pe ucenicul pe care-l iubea stând acolo, a zis mamei Sale:
{\WJ{“Femeie, iată fiul tău!”
}}

\VS{27}Apoi a zis ucenicului:
{\WJ{“Iată mama ta!”. Din acel ceas, discipolul a luat-o la el acasă.
}}


\par }{\PP \VS{28}După aceea, Isus, văzând că toate lucrurile se terminaseră, ca să se împlinească Scriptura, a zis:
{\WJ{“Mi-e sete!”
}}

\VS{29}Și era pus acolo un vas plin cu oțet; au pus un burete plin de oțet pe isop și l-au ținut la gura Lui.

\VS{30}După ce a primit deci oțetul, Isus a zis: “S-a
{\WJ{sfârșit!” Apoi și-a plecat capul și și-a dat duhul.
}}


\par }{\PP \VS{31}De aceea iudeii, pentru că era ziua pregătirii, ca să nu rămână trupurile pe cruce în Sabat, căci Sabatul acela era un Sabat special, au cerut lui Pilat să li se rupă picioarele și să fie luați.

\VS{32}Soldații au venit, așadar, și au rupt picioarele celui dintâi și ale celuilalt care fusese răstignit împreună cu el;

\VS{33}dar, când au ajuns la Isus și au văzut că era deja mort, nu i-au rupt picioarele.

\VS{34}Totuși, unul dintre soldați i-a străpuns coasta cu o suliță și imediat a ieșit sânge și apă.

\VS{35}Cel care a văzut a depus mărturie, și mărturia lui este adevărată. El știe că spune adevărul, ca să credeți.

\VS{36}Căci aceste lucruri s-au întâmplat pentru ca să se împlinească Scriptura: “Nici un os din el nu va fi zdrobit”.

\VS{37}Și iarăși o altă Scriptură spune: “Se vor uita la Cel pe care L-au străpuns.”


\par }{\PP \VS{38}După acestea, Iosif din Arimateea, care era ucenic al lui Isus, dar care, de frica iudeilor, a cerut lui Pilat să ia trupul lui Isus. Pilat i-a dat permisiunea. A venit, așadar, și a luat trupul Lui.

\VS{39}Nicodim, care la început venise la Isus noaptea, a venit și el aducând un amestec de smirnă și aloe, cam o sută de lire romane.

\VS{40}Au luat deci trupul lui Isus și l-au legat în pânze de in cu mirodenii, așa cum este obiceiul iudeilor de a-l îngropa.

\VS{41}În locul unde a fost răstignit era o grădină. În grădină era un mormânt nou, în care nu mai fusese așezat niciun om până atunci.

\VS{42}Atunci, din pricina zilei de pregătire a iudeilor (căci mormântul era aproape), au pus pe Isus acolo.



\par }\Chap{20}{\PP \VerseOne{1}În prima zi a săptămânii, Maria Magdalena s-a dus dis-de-dimineață la mormânt, pe când era încă întuneric, și a văzut că piatra fusese scoasă de pe mormânt.

\VS{2}De aceea a alergat și a venit la Simon Petru și la celălalt ucenic pe care-l iubea Isus și le-a zis: “L-au luat pe Domnul din mormânt și nu știm unde L-au pus!”


\par }{\PP \VS{3}Atunci Petru și celălalt ucenic au ieșit și s-au dus la mormânt.

\VS{4}Și amândoi alergau împreună. Celălalt ucenic l-a întrecut pe Petru și a ajuns primul la mormânt.

\VS{5}S-a aplecat și a privit înăuntru, a văzut pânzele de in așezate acolo, dar n-a intrat.

\VS{6}Atunci Simon Petru a venit după el și a intrat în mormânt. A văzut așezate pânzele de in,

\VS{7}iar pânza care fusese pe capul Lui, nu era așezată împreună cu pânzele de in, ci înfășurată într-un loc de sine stătător.

\VS{8}Atunci a intrat și celălalt ucenic care venise primul la mormânt; a văzut și el și a crezut.

\VS{9}Căci încă nu cunoșteau Scriptura, că trebuie să învieze din morți.

\VS{10}Ucenicii s-au dus din nou la casele lor.


\par }{\PP \VS{11}Dar Maria stătea afară la mormânt și plângea. Și, în timp ce plângea, s-a aplecat și s-a uitat în mormânt,

\VS{12}și a văzut doi îngeri în alb care stăteau, unul la cap și altul la picioare, unde fusese așezat trupul lui Isus.

\VS{13}Ei au întrebat-o: “Femeie, de ce plângi?”.

\par }{\PP Ea le-a spus: “Pentru că l-au luat pe Domnul meu și nu știu unde l-au pus”.

\VS{14}După ce a spus aceasta, s-a întors și a văzut pe Isus în picioare, și nu știa că este Isus.


\par }{\PP \VS{15}Isus i-a zis: “{\WJ{Femeie, de ce plângi? Pe cine cauți?”
}}

\par }{\PP Ea, presupunând că este grădinarul, i-a spus: “Domnule, dacă l-ai dus, spune-mi unde l-ai pus și îl voi lua.”


\par }{\PP \VS{16}Isus i-a zis:
{\WJ{“Maria.”
}}

\par }{\PP Ea s-a întors și i-a spus: “Rabboni!”, adică “Învățătorule!”.


\par }{\PP \VS{17}Isus i-a zis:
{\WJ{“Nu Mă țineți, căci încă nu M-am înălțat la Tatăl Meu, ci du-te la frații Mei și spune-le: Mă înalț la Tatăl Meu și Tatăl vostru, la Dumnezeul Meu și Dumnezeul vostru.”
}}


\par }{\PP \VS{18}Maria Magdalena a venit și a spus ucenicilor că a văzut pe Domnul și că El i-a spus aceste lucruri.

\VS{19}Așadar, când s-a făcut seară în ziua aceea, prima zi a săptămânii, și când ușile erau încuiate acolo unde erau adunați discipolii, de frica iudeilor, Isus a venit, a stat în mijloc și le-a zis:
{\WJ{“Pace vouă!”.
}}


\par }{\PP \VS{20}După ce a zis acestea, le-a arătat mâinile și coasta Lui. Ucenicii s-au bucurat, așadar, când L-au văzut pe Domnul.

\VS{21}Isus le-a zis din nou: “{\WJ{Pace vouă! Cum M-a trimis Tatăl, așa vă trimit și Eu pe voi”.
}}

\VS{22}După ce a spus acestea, a suflat peste ei și le-a zis:
{\WJ{“Primiți Duhul Sfânt!
}}

\VS{23}{\WJ{Dacă iertați păcatele cuiva, i-au fost iertate. Dacă rețineți păcatele cuiva, ele au fost reținute.”
}}


\par }{\PP \VS{24}Dar Toma, unul din cei doisprezece, numit Didimus, nu era cu ei când a venit Isus.

\VS{25}De aceea, ceilalți discipoli i-au spus: “L-am văzut pe Domnul!”

\par }{\PP Dar el le-a zis: “Dacă nu voi vedea în mâinile Lui amprenta cuielor, dacă nu voi pune degetul meu în amprenta cuielor și dacă nu voi pune mâna mea în coasta Lui, nu voi crede.”


\par }{\PP \VS{26}După opt zile, ucenicii Lui au intrat din nou înăuntru și Toma era cu ei. Isus a venit, ușile fiind încuiate, a stat în mijloc și a zis:
{\WJ{“Pace vouă!”.
}}

\VS{27}Apoi i-a spus lui Toma: “{\WJ{Întinde aici degetul tău și vezi mâinile Mele. Întinde aici mâna ta și pune-o în coasta Mea. Nu fi necredincios, ci credincios.”
}}


\par }{\PP \VS{28}Toma I-a răspuns: “Domnul meu și Dumnezeul meu!”


\par }{\PP \VS{29}Isus i-a zis: “{\WJ{Pentru că M-ai văzut, ai crezut. Ferice de cei care nu au văzut și au crezut.”
}}


\par }{\PP \VS{30}Isus a mai făcut multe alte semne în fața ucenicilor Săi, care nu sunt scrise în cartea aceasta;

\VS{31}dar acestea sunt scrise pentru ca voi să credeți că Isus este Hristosul, Fiul lui Dumnezeu, și, crezând, să aveți viață în Numele Lui.



\par }\Chap{21}{\PP \VerseOne{1}După acestea, Isus S-a arătat din nou ucenicilor la Marea Tiberiadei. El s-a revelat în felul următor.

\VS{2}Simon Petru, Toma zis Didim, Natanael din Cana Galileii, fiii lui Zebedei și alți doi dintre discipolii Săi erau împreună.

\VS{3}Simon Petru le-a spus: “Mă duc la pescuit”.

\par }{\PP I-au spus: “Și noi venim cu tine”. Au ieșit imediat și au intrat în barcă. În acea noapte, nu au prins nimic.

\VS{4}Dar când deja se făcuse ziuă, Isus stătea pe plajă; totuși, discipolii nu știau că era Isus.

\VS{5}Isus le-a zis deci:
{\WJ{“Copii, aveți ceva de mâncare?”
}}

\par }{\PP Ei i-au răspuns: “Nu”.


\par }{\PP \VS{6}Și le-a zis: “{\WJ{Aruncați năvodul în partea dreaptă a corăbiei și veți găsi câteva.”
}}

\par }{\PP L-au aruncat, deci, și acum nu mai puteau să-l atragă pentru mulțimea de pești.

\VS{7}De aceea, ucenicul acela pe care-l iubea Isus i-a spus lui Petru: “Este Domnul!”

\par }{\PP Când Simon Petru a auzit că este Domnul, și-a înfășurat haina (căci era gol) și s-a aruncat în mare.

\VS{8}Dar ceilalți ucenici au venit cu barca mică (căci nu erau departe de uscat, ci la vreo două sute de coți), trăgând plasa plină de pește.

\VS{9}Când au ieșit pe uscat, au văzut acolo un foc de cărbuni, pe care erau așezate pește și pâine.

\VS{10}Isus le-a zis:
{\WJ{“Aduceți o parte din peștele pe care tocmai l-ați prins.”
}}


\par }{\PP \VS{11}Simon Petru s-a suit și a tras la țărm o plasă plină cu o sută cincizeci și trei de pești mari. Chiar dacă erau atât de mulți, plasa nu s-a rupt.


\par }{\PP \VS{12}Isus le-a zis: “{\WJ{Veniți și luați micul dejun!”
}}

\par }{\PP Nici unul dintre ucenici nu a îndrăznit să-l întrebe: “Cine ești tu?”, știind că era Domnul.


\par }{\PP \VS{13}Isus a venit, a luat pâinea, le-a dat-o și peștele, și tot așa și peștele.

\VS{14}Aceasta este acum a treia oară când Iisus se arată discipolilor Săi după ce a înviat din morți.

\VS{15}După ce au luat micul dejun, Isus i-a zis lui Simon Petru:
{\WJ{“Simon, fiul lui Iona, mă iubești tu mai mult decât aceștia?”
}}

\par }{\PP El i-a răspuns: “Da, Doamne, știi că am o afecțiune pentru tine”.

\par }{\PP El i-a spus:
{\WJ{“Paște mieii mei”.
}}

\VS{16}Și i-a zis a doua oară:
{\WJ{“Simon, fiul lui Iona, mă iubești?”
}}

\par }{\PP El i-a răspuns: “Da, Doamne, știi că am o afecțiune pentru tine”.

\par }{\PP El i-a spus:
{\WJ{“Păzește oile mele”.
}}

\VS{17}A treia oară i-a zis:
{\WJ{“Simon, fiul lui Iona, mă iubești?”
}}

\par }{\PP Petru a fost mâhnit pentru că l-a întrebat a treia oară:
{\WJ{“Mă iubești?”. El i-a răspuns: “Doamne, Tu știi totul. Tu știi că am afecțiune pentru tine”.
}}

\par }{\PP Isus i-a spus:
{\WJ{“Paște oile Mele.
}}

\VS{18}{\WJ{Adevărat îți spun că, atunci când erai tânăr, te îmbrăcai singur și mergeai pe unde voiai. Dar când vei fi bătrân, vei întinde mâinile și altul te va îmbrăca și te va duce unde nu vrei să mergi.”
}}


\par }{\PP \VS{19}Și a spus aceasta, ca să arate prin ce fel de moarte vrea să slăvească pe Dumnezeu. După ce a spus aceasta, i-a zis:
{\WJ{“Urmează-mă”.
}}


\par }{\PP \VS{20}Atunci Petru, întorcându-se, a văzut un ucenic care îl urmărea. Acesta era ucenicul pe care Isus îl iubea, cel care se aplecase și el la pieptul lui Isus la cină și care întrebase: “Doamne, cine te va trăda?”

\VS{21}Petru, văzându-l, l-a întrebat pe Isus: “Doamne, ce-i cu omul acesta?”


\par }{\PP \VS{22}Isus i-a zis: “{\WJ{Dacă vreau să rămână până ce voi veni Eu, ce-ți pasă ție? Urmează-mă tu!”
}}

\VS{23}De aceea s-a răspândit acest cuvânt printre frați, ca acest ucenic să nu moară. Totuși, Isus nu i-a spus că nu va muri, ci:
{\WJ{“Dacă vreau ca el să rămână până când voi veni eu, ce-ți pasă ție?”
}}


\par }{\PP \VS{24}Acesta este ucenicul care mărturisește despre aceste lucruri și care a scris aceste lucruri. Noi știm că mărturia lui este adevărată.

\VS{25}Mai sunt și multe alte lucruri pe care le-a făcut Isus, care, dacă ar fi scrise toate, cred că nici măcar lumea însăși nu ar avea loc pentru cărțile care s-ar scrie.


\par }